\documentclass[11pt]{article}

\usepackage[T1]{fontenc}
\usepackage[utf8]{inputenc}
\usepackage[english]{babel}
\usepackage[margin=1in]{geometry}
\usepackage{lmodern}
\usepackage{microtype}
\usepackage{mathtools}
\usepackage{amsmath,amssymb,amsthm}
\usepackage{array}
\usepackage{booktabs}
\usepackage{tabularx}
\usepackage{float}
\usepackage{graphicx}
\usepackage{pdflscape}
\usepackage{caption}
\usepackage{enumitem}
\usepackage{fancyhdr}
\usepackage[numbers]{natbib}
\usepackage{doi}
\usepackage{xurl}
\usepackage{hyperref}
\usepackage{cleveref}
\usepackage{orcidlink}
\usepackage{titling}



\DeclareMathOperator{\Capacity}{Cap}
\DeclareMathOperator{\Cost}{Cost}
\DeclareMathOperator{\Ran}{Ran}
\DeclareMathOperator{\Ker}{Ker}
\DeclareMathOperator{\tr}{tr}
\DeclareMathOperator{\dist}{dist}

\newcommand{\K}{\mathsf K}
\newcommand{\XiRes}{\Xi}
\newcommand{\preceqq}{\preceq}
\newcommand{\eps}{\varepsilon}

\newcommand{\PreSB}{\mathbf{PreSB}}
\newcommand{\AOR}{\mathbf{AOR}}
\newcommand{\RefPreSB}{\mathbf{RefPreSB}}
\newcommand{\CascadeStableAOR}{\mathbf{AOR}^{\mathrm{casc}}}
\newcommand{\RefStableAOR}{\mathbf{AOR}^{\mathrm{ref}}}

\newcommand{\ClSB}{\mathrm{Cl}_{SB}}

\newcommand{\XiSB}{\Xi_{SB}}
\newcommand{\Sigobs}{\Sigma_{\mathrm{obs}}}
\newcommand{\Sigstage}{\Sigma_{\mathrm{stage}}}
\newcommand{\Sigconstr}{\Sigma_{\mathrm{constraint}}}
\newcommand{\Sigrewrite}{\Sigma_{\mathrm{rewrite}}}
\newcommand{\Sighol}{\Sigma_{\mathrm{holonomy}}}
\newcommand{\Sigres}{\Sigma_{\mathrm{residual}}}
\newcommand{\Signon}{\Sigma_{\mathrm{nonclaim}}}
\newcommand{\Sigmeta}{\Sigma_{\mathrm{meta}}}

\newcommand{\Ddescent}{\Delta_{\mathrm{descent}}}
\newcommand{\Dtransport}{\Delta_{\mathrm{transport}}}
\newcommand{\Dinterface}{\Delta_{\mathrm{interface}}}
\newcommand{\Dlocal}{\Delta_{\mathrm{local}}}
\newcommand{\Dglobal}{\Delta_{\mathrm{global}}}
\newcommand{\Dpresentation}{\Delta_{\mathrm{presentation}}}
\newcommand{\Dsource}{\Delta_{\mathrm{source}}}
\newcommand{\Dhorizon}{\Delta_{\mathrm{horizon}}}
\newcommand{\Dscale}{\Delta_{\mathrm{scale}}}
\newcommand{\Dlimit}{\Delta_{\mathrm{limit}}}
\newcommand{\Dsym}{\Delta_{\mathrm{sym}}}
\newcommand{\Drole}{\Delta_{\mathrm{role}}}
\newcommand{\Dtarget}{\Delta_{\mathrm{target}}}

\newcommand{\zeroSt}{\texttt{zero}}
\newcommand{\budgeted}{\texttt{budgeted}}
\newcommand{\bridged}{\texttt{bridged}}
\newcommand{\hidden}{\texttt{hidden}}
\newcommand{\moduli}{\texttt{moduli}}
\newcommand{\selected}{\texttt{selected}}
\newcommand{\probabilistic}{\texttt{probabilistic}}
\newcommand{\outOfScope}{\texttt{out\_of\_scope}}
\newcommand{\blocked}{\texttt{blocked}}
\newcommand{\approvedOther}{\texttt{approved\_other}}

\newcommand{\asymptoticBudgeted}{\texttt{asymptotic\_budgeted}}
\newcommand{\summable}{\texttt{summable}}
\newcommand{\contractive}{\texttt{contractive}}
\newcommand{\eventualZero}{\texttt{eventual\_zero}}
 
\newtheorem{theorem}{Theorem}[section]
\newtheorem{lemma}[theorem]{Lemma}
\newtheorem{proposition}[theorem]{Proposition}
\newtheorem{corollary}[theorem]{Corollary}
\theoremstyle{definition}
\newtheorem{definition}[theorem]{Definition}
\newtheorem{remark}[theorem]{Remark}
\newtheorem{nonclaim}[theorem]{Nonclaim}
\floatstyle{ruled}
\newfloat{algorithm}{tbp}{loa}
\floatname{algorithm}{Algorithm}

\urlstyle{same}
\captionsetup{
  font=small,
  labelfont=bf,
  labelsep=period
}
\captionsetup[table]{skip=6pt,position=top}
\captionsetup[figure]{skip=6pt}
\captionsetup[algorithm]{skip=6pt,position=top}
\crefname{algorithm}{algorithm}{algorithms}
\Crefname{algorithm}{Algorithm}{Algorithms}
\crefname{theorem}{theorem}{theorems}
\Crefname{theorem}{Theorem}{Theorems}
\crefname{lemma}{lemma}{lemmas}
\Crefname{lemma}{Lemma}{Lemmas}
\crefname{proposition}{proposition}{propositions}
\Crefname{proposition}{Proposition}{Propositions}
\crefname{corollary}{corollary}{corollaries}
\Crefname{corollary}{Corollary}{Corollaries}
\crefname{definition}{definition}{definitions}
\Crefname{definition}{Definition}{Definitions}
\crefname{remark}{remark}{remarks}
\Crefname{remark}{Remark}{Remarks}
\setlist[enumerate]{leftmargin=*, itemsep=2pt, topsep=4pt}
\setlength{\droptitle}{-3.2em}
\pretitle{\begin{center}\LARGE\bfseries}
\posttitle{\par\end{center}\vspace{0.5em}}
\preauthor{\begin{center}\normalsize}
\postauthor{\par\end{center}\vspace{-0.4em}}
\predate{\begin{center}\small}
\postdate{\par\end{center}\vspace{0.6em}}
\setlength{\emergencystretch}{2em}
\clubpenalty=10000
\widowpenalty=10000
\displaywidowpenalty=10000
\interfootnotelinepenalty=10000
\allowdisplaybreaks[2]
\fancypagestyle{firstpage}{\fancyhf{}
  \fancyfoot[C]{\scriptsize Automorph Inc., Wilmington, DE, USA \quad \texttt{ioannis@automorph.io} \quad \textcopyright\ 2026 \quad CC-BY 4.0}
  \renewcommand{\headrulewidth}{0pt}
  \renewcommand{\footrulewidth}{0pt}
}

\title{Audited Operational Realisability:\\
A Closure-Completion Reflection for Carriers with Audit Data}
\author{Ioannis Tsiokos\,\orcidlink{0009-0009-7659-5964}}
\date{\today}
\hypersetup{
  pdftitle={Audited Operational Realisability: A Closure-Completion Reflection for Carriers with Audit Data},
  pdfauthor={Ioannis Tsiokos},
  hidelinks
}

\begin{document}

\maketitle
\thispagestyle{firstpage}

\begin{abstract}
Many mathematical frameworks come with audit data alongside their
main objects: source ledgers, route registries, residual budgets,
status annotations, nonclaim boundaries, and meta-audit records.
\emph{Audited Operational Realisability} (AOR) is the meta-theory of
typed carriers that close as audit systems within a declared scope.
A \emph{PreSB carrier} is a typed carrier together with this audit
data; it is an \emph{AOR carrier} when every declared audit
obligation is discharged inside the scope.  We show that every PreSB
carrier \(\PreSB_{\mathfrak S}\) admits a canonical completion into
the full subcategory of AOR carriers via a reflection
\(\ClSB\dashv\iota\) obtained as a monotone Knaster--Tarski
saturation on the complete lattice of AOR-extending operations.  The
reflection is developed in four parts: the closure-completion
construction and its universal property; the eight-stratum
decomposition of the structural defect \(\XiSB\); the twelve-type
\(\Sigma_{\mathrm{residual}}\) atlas with its discharge cascades and
cascade-stable AOR; and refinement chains with asymptotic statuses,
yielding the hierarchy
\(\RefStableAOR\subseteq\CascadeStableAOR\subseteq\AOR\).  The
development is meta-theoretic and carries no PDE-specific worked
example.  An AOR-instance presentation of the Navier--Stokes
smooth-data carrier is given in \citet{Tsiokos2026NavierStokes},
which cites the present meta-theory at paper level.  Two scope
choices apply.  Global cascade confluence is established by Newman's
lemma applied to the local diamond property and the well-founded
discharge relation, both of which we prove explicitly.  The hierarchy
is exhibited via named witness carriers and their typed obstruction
markers; full predicate-level non-membership proofs at every level
remain framework future work.
\end{abstract}

\medskip
\noindent\textbf{Keywords.} Six Birds theory; audited operational realisability; PreSB; closure completion; Knaster--Tarski reflection; residual atlas; discharge cascade; refinement-stable AOR.

\section{Introduction}\label{sec:aor-intro}

Fix a scope record \(\mathfrak S\) declaring the admissible
carrier-class, source ledger, interface family, and route registry.
The category \(\PreSB_{\mathfrak S}\) consists of typed carriers
equipped with declared audit data over \(\mathfrak S\) --- observation,
stage, constraint, rewrite, holonomy, residual, status, nonclaim, and
meta-audit registries --- together with morphisms preserving the
scoped package.  Its full subcategory \(\AOR_{\mathfrak S}\) contains
the audit-closed carriers, those whose declared records, routes,
residuals, nonclaims, and meta-audits have all been typed, registered,
statused, or discharged inside \(\mathfrak S\).  The central question
is the corresponding closure problem: when does an arbitrary object of
\(\PreSB_{\mathfrak S}\) admit a canonical completion into
\(\AOR_{\mathfrak S}\), and what is the structure of that completion?

The construction is a categorical reflection.  The functor \(\ClSB\)
sends a PreSB carrier to its closure-completed AOR carrier, and the
inclusion \(\iota:\AOR_{\mathfrak S}\hookrightarrow
\PreSB_{\mathfrak S}\) is its right adjoint, so that maps from a
PreSB carrier into an audit-closed carrier factor uniquely through the
closure-completed object.  What is specific to the present framework
is the audit content of that closure: the completion is not a formal
hull but the least saturation that adds the AOR-extending operations
required by the carrier's records, residuals, and status obligations.

The first diagnostic device is the structural defect \(\XiSB\).  For a
PreSB carrier, \(\XiSB\) measures the aggregate failure of audit
closure, split into eight strata: observation, stage, constraint,
rewrite, holonomy, residual, nonclaim, and meta-audit defects.  AOR
membership therefore decomposes: it holds exactly when every stratum
is empty.  The residual stratum is refined further by a twelve-type
residual atlas, which separates failures of descent, transport,
interface, local and global consistency, presentation, source
admission, horizon, scale, limit, symmetry, role, and target typing.
Residual atoms carry statuses drawn from the AOR-8 catalogue and may
generate downstream obligations, so audit closure is also an
operational question about how discharge propagates.

The second operational layer is the discharge cascade.  A cost record
\(\Cost(t,s)\) assigns to a residual type \(t\) and status \(s\) the
finite multiset of downstream obligations produced by that discharge
choice; the operator \(\mathrm{Down}\) collects these obligations
across a witness family.  The resulting cascade graph is bipartite,
alternating residual atoms with discharge witnesses.  Cascade-stable
AOR, \(\CascadeStableAOR\), consists of those audit-closed carriers
that come equipped with a declared discharge assignment whose cascades
terminate under the cost discipline.  The status assignment is part
of the data of the carrier; it cannot be chosen after the fact to
secure cascade stability.

The third layer is asymptotic.  Refinement morphisms annotate
scope-preserving PreSB morphisms by one of six refinement kinds, and
directed refinement chains form the category \(\RefPreSB_{\mathfrak S}\).
At directed colimits, finite cascade closure is not always enough; limit
residuals require one of the asymptotic statuses
\(\asymptoticBudgeted\), \(\summable\), \(\contractive\), or
\(\eventualZero\).  The corresponding subclass
\(\RefStableAOR\) consists of carriers whose audit closure survives both
finite cascade discharge and refinement-colimit passage.  The resulting
hierarchy is
\[
  \RefStableAOR_{\mathfrak S}
  \subseteq
  \CascadeStableAOR_{\mathfrak S}
  \subseteq
  \AOR_{\mathfrak S},
\]
and the witness carriers of \Cref{thm:aor:hierarchy} mark the two
boundaries between ordinary audit closure, finite cascade closure, and
refinement-stable closure.

The motivation for the abstraction comes from earlier audited
frameworks in the same programme.  The formed-layer membrane theory
of \citet{Tsiokos2026NeedleKiller} and the adequacy-residual calculus
of \citet{Tsiokos2026XiCompanion} both rely on a disciplined account
of what has been registered, what has been transported, what residual
remains, and what is explicitly out of scope.  The present paper
abstracts that bookkeeping into a self-contained category-theoretic
and order-theoretic meta-theory.  It does not reprove results from
those works, nor does any theorem here depend on their internal
labels.  In the opposite direction, \citet{Tsiokos2026NavierStokes}
uses the AOR meta-theory as a consumer framework for a presentation
of the Navier--Stokes smooth-data carrier as an AOR instance; the
PDE-specific closure stack is developed there, not here.

One notational ambiguity is worth flagging at the outset.  The symbol
\(\XiSB\) in this paper denotes the eight-stratum structural defect of
a PreSB carrier; it is not the adequacy residual \(\Xi_C(D\mid L)\) of
\citet{Tsiokos2026XiCompanion}.  The two objects share the Greek
letter because both measure a kind of leftover obligation, but they
live in disjoint formal contexts: \(\XiSB\) belongs to the AOR
meta-theory of typed audit records, while \(\Xi_C(D\mid L)\) belongs
to the Schur-complement calculus of native and dissolving probe
families.

A formal companion of the development is recorded in
\Cref{app:aor-formalization}; the body is self-contained and does not
depend on it.

The remainder of the paper proceeds as follows.  \Cref{sec:aor-background}
sets up the mathematical entry point: scoped carriers, reflective
completion, least fixed points, well-founded discharge, and directed
colimits.  \Cref{sec:aor-problem-setup} defines PreSB carriers, their
morphisms, and the AOR full subcategory.  \Cref{sec:aor-methods}
constructs \(\ClSB\), states its universal property, and packages the
adjunction with the inclusion.  \Cref{sec:aor-primary-results} develops the
primary diagnostic layer: \(\XiSB\), the eight strata, the residual atlas,
the atom calculus, and the AOR-8 status catalogue.  \Cref{sec:aor-secondary-results}
builds discharge cascades and cascade-stable AOR.  \Cref{sec:aor-analysis}
adds refinement chains, asymptotic statuses, directed-colimit behavior, and
the hierarchy.  Finally, \Cref{sec:aor-discussion} records the status,
scope, and nonclaim boundary, and \Cref{sec:aor-conclusion} summarizes the
audit-closure meta-theory.
 \section{Background}\label{sec:aor-background}

The framework developed here abstracts an audit discipline that
already appears in the formed-layer membrane theory of
\citet{Tsiokos2026NeedleKiller}.  In that setting the mathematical
object of interest is never just a bare carrier; it is accompanied by
records of which probes are native, which routes are legal, which
residuals remain unpaid, which statements are explicitly outside
scope, and which bookkeeping obligations have been accepted.  We do
not rederive that membrane theory, but isolate the meta-theoretic
pattern behind it: a carrier can be meaningfully used only after its
record layer has been audited inside a declared scope.

The adequacy-residual calculus of \citet{Tsiokos2026XiCompanion}
supplies a second motivating source.  There the central object is the
Schur-complement residual \(\Xi_C(D\mid L)\), measuring what a
dissolving probe family \(D\) cannot explain through a native probe
family \(L\) under an audit operator \(C\).  The calculus distinguishes
native explanation from leftover residual, exact adequacy from
defective adequacy, and lawful promotion from blind-spot failure.
AOR neither imports those theorems nor depends on them; it abstracts
the surrounding bookkeeping question: once a residual, route, status,
or nonclaim has been declared, what does it mean for the whole
carrier to be audit-closed?

The categorical language used below is standard
\citep{MacLane1998Categories,Riehl2016CategoryTheory,Borceux1994CategoricalAlgebra}.
A full subcategory keeps all morphisms between the selected objects; in
this paper the selected objects are the audit-closed carriers.  A
reflective subcategory is a full subcategory whose inclusion admits a
left adjoint, so every ambient object has a universal approximation
inside the subcategory.  The reflector here is the closure-completion
functor \(\ClSB\), and the inclusion is
\(\iota:\AOR_{\mathfrak S}\hookrightarrow\PreSB_{\mathfrak S}\).  The
identity-field counit says that an already audit-closed carrier is
unchanged by applying the closure completion.  Directed colimits enter
later when a carrier is not built in one step but as the limit of a
refinement chain; the relevant background is the theory of accessible
and locally presentable categories \citep{AdamekRosicky1994LocallyPresentable}.

The order-theoretic input is equally classical
\citep{DaveyPriestley2002Lattices,GierzEtAl2003ContinuousLattices}.  A
complete lattice is a partially ordered set in which every family has a
meet and a join.  A monotone endomap on such a lattice has least and
greatest fixed points by the Knaster--Tarski theorem
\citep{Knaster1928SetFunctions,TarskiFixedPoint1955}.  The construction of
\(\ClSB\) uses this pattern in the form of monotone saturation:
starting from a seed carrier, one repeatedly adds the minimal
AOR-extending operations forced by the audit data, and the least fixed
point is the closure-completed carrier.  The constructive ordinal
iteration giving this fixed point is the version of
\citet{CousotCousotConstructive1979}; the same lattice-fixed-point
machinery underlies abstract interpretation
\citep{CousotCousot1977AbstractInterpretation}.  The body uses the
fixed-point theorem as standard order theory.

The defect layer should be read as typed diagnostic bookkeeping
rather than as an analytic norm or an energy.  The structural defect
\(\XiSB\) is a tagged aggregate of audit failures, and its eight
strata separate observation typing, stage data, constraints, rewrites,
holonomy records, residuals, nonclaims, and meta-audits.  This
separation prevents a single error class from hiding several distinct
audit obligations, and it makes the phrase ``defect zero'' operational:
\(\XiSB=0\) means that each stratum has independently emptied.

Residual discharge is modeled by a finite labelled rewriting discipline.
A residual atom receives a status, a status produces downstream obligations
through the cost table, and the downstream operator collects the next layer
of atoms.  This is why the cascade graph is bipartite: its layers alternate
between residual atoms and discharge witnesses.  Termination is expressed by
a well-founded measure on the residual type and status data.  The relevant
background is the standard rewriting-theory package in which termination
and local confluence control normal-form behavior; see
\citet{KnuthBendix1970SimpleWord} and \citet{Huet1980ConfluentReductions},
the modern treatments in \citet{BaaderNipkow1998TermRewriting} and
\citet{Terese2003TermRewriting}, and Newman's lemma in its original form
\citep{Newman1942Combinatorial}.

Refinement adds one more layer.  A refinement morphism is a
scope-preserving map of carriers annotated by the kind of refinement being
performed: expanding the scope, adding records, growing stages, tightening
constraints, strengthening morphisms, or refining the carrier data.  A
refinement chain is a directed diagram of such morphisms.  At a directed
colimit, record data may be obtained by directed union while carrier data
may be represented by compatible threads.  Finite cascade discharge need not
survive this passage automatically, because a residual can appear only at
the tail or limit stage.  The asymptotic statuses
\(\asymptoticBudgeted\), \(\summable\), \(\contractive\), and
\(\eventualZero\) are the vocabulary used to pay these limit-stage
obligations.

These ingredients yield the hierarchy developed later:
\[
  \RefStableAOR_{\mathfrak S}
  \subseteq
  \CascadeStableAOR_{\mathfrak S}
  \subseteq
  \AOR_{\mathfrak S}.
\]
Ordinary AOR membership means zero structural defect.  Cascade-stable
AOR adds a declared discharge assignment whose finite cascades close.
Refinement-stable AOR further requires compatibility with refinement
chains and their asymptotic residuals.  The witness boundaries for these
inclusions are exhibited in \Cref{sec:aor-analysis} via named carriers.

The relationship to \citet{Tsiokos2026NavierStokes} is one-way.  That
paper consumes the AOR vocabulary to present the Navier--Stokes
smooth-data carrier as a refinement-stable AOR instance.  No
theorem of the present paper depends on that presentation, and no
PDE-specific worked example is carried here.
 \section{Problem setup: PreSB carriers and the AOR full subcategory}\label{sec:aor-problem-setup}

\Cref{sec:aor-background} described audit closure in ordinary
mathematical terms.  This section fixes the formal objects used in
the rest of the paper.  A PreSB carrier is not only a carrier space
with quotient data; it is a scoped record package carrying the
registries, residuals, statuses, nonclaims, and meta-audit entries
that must be checked before the carrier can be treated as
audit-closed.  The scope record \(\mathfrak S\) is part of the data of
the carrier rather than an external side condition: two carriers
whose admissible sources, interfaces, routes, or carrier classes are
declared in different scopes are not compared inside a common
category.

The carrier data come first.  The next definition lists the seventeen
fields that the paper audits.  The symbols \(H,q,Q\) encode the
legal-quotient carrier data; \(\Gamma\) is the route registry; and
the source, interface, observation, stage, constraint, rewrite,
holonomy, residual, status, nonclaim, and meta-audit registries form
the record layer to which the closure-completion construction will be
applied.

\begin{definition}[PreSB carrier]\label{def:aor:presb-carrier}
A \emph{PreSB carrier} is a tuple
\[
  \mathsf C = (H,\,q,\,Q,\,\Gamma,\,\mathsf R,\,\mathsf C\text{-reg},\,
                \mathrm{Src},\,\mathrm{Ifc},\,\mathrm{Obs},\,
                \mathrm{Stage},\,\mathrm{Constr},\,\mathrm{Rewrite},\,
                \mathrm{Hol},\,\mathrm{Res},\,\mathrm{Stat},\,
                \mathrm{NC},\,\mathrm{Meta};\,\mathfrak S),
\]
of seventeen typed fields: legal-quotient carrier and audit data
$(H,q,Q)$; route registry $\Gamma$; record registry $\mathsf R$ and
claim registry $\mathsf C$-reg; source, interface, and observation
registries $(\mathrm{Src},\mathrm{Ifc},\mathrm{Obs})$; stage, constraint,
rewrite, and holonomy registries
$(\mathrm{Stage},\mathrm{Constr},\mathrm{Rewrite},\mathrm{Hol})$;
residual and status registries $(\mathrm{Res},\mathrm{Stat})$;
nonclaim register $\mathrm{NC}$ and meta-audit register
$\mathrm{Meta}$; together with the scope record $\mathfrak S$ on which
the carrier is evaluated.
\end{definition}

The scope record is named separately because it controls the
formation of both objects and morphisms.  A carrier-class declaration,
source ledger, interface family, and route registry acquire their
intended meaning only after the scope has fixed which such data are
admissible.

\begin{definition}[Scope record]\label{def:aor:scope-record}
A \emph{scope record} $\mathfrak S$ is a typed declaration of the
admissible carrier-class, the admissible source ledger, the admissible
interface family, and the admissible route registry against which a
PreSB carrier is audited.  Scope is a first-class field of a PreSB
carrier; comparing two carriers of incompatible scopes is undefined.
\end{definition}

Morphisms preserve this scoped audit package.  They are not arbitrary maps
between carriers, but tuples of field maps that respect the quotient,
transport, and audit-pairing structure already declared in the objects.
Thus, after a scope has been fixed, the PreSB carriers form an ordinary
category with componentwise identities and composition.

\begin{definition}[PreSB morphism]\label{def:aor:morphism-presb}
A \emph{PreSB morphism} $f:\mathsf C\to\mathsf C'$ between two PreSB
carriers of identical scope is a tuple of typed scope-preserving
field morphisms that commute with the carrier quotient maps, the route
registry transports, and the audit form pairings.  Morphisms compose
componentwise; the identity morphism is the tuple of identity
field-maps.  The category $\PreSB_{\mathfrak S}$ has PreSB carriers of
scope $\mathfrak S$ as objects and the above morphisms as arrows.
\end{definition}

The audit-closed objects are then selected by a defect-zero
condition.  The definition states this condition through the
structural defect, whose eight strata are developed in
\Cref{sec:aor-primary-results-defect}.  At this point the equivalent
operational reading is the more useful one: the declared records are
typed, the cited routes are audited, the residuals have statuses, the
nonclaims have entries, and the meta-audit obligations have been
discharged.

\begin{definition}[AOR subcategory]\label{def:aor:aor-subcategory}
The \emph{AOR subcategory} $\AOR_{\mathfrak S}$ is the full
subcategory of $\PreSB_{\mathfrak S}$ whose objects are the PreSB
carriers $\mathsf C$ for which the eight-stratum structural defect
$\XiSB(\mathsf C)$ (\Cref{def:aor:xi-sb}) is zero.  Equivalently, every
declared record of $\mathsf C$ is typed, every cited route is audited,
every residual is statused, every nonclaim is registered, and every
meta-audit obligation is discharged.
\end{definition}

Because \(\AOR_{\mathfrak S}\) is full, it inherits the ambient
morphisms between selected objects.  What remains to verify is
isomorphism-invariance of membership: if two scoped carriers are
isomorphic in \(\PreSB_{\mathfrak S}\), then the typed audit
obligations transported along the isomorphism vanish on one carrier
exactly when they vanish on the other.

\begin{theorem}[AOR is a full subcategory]\label{thm:aor:aor-full-subcat}
The inclusion $\iota:\AOR_{\mathfrak S}\hookrightarrow\PreSB_{\mathfrak S}$
is a fully faithful functor.  The defect predicate $\XiSB(\mathsf C)=0$
is preserved by isomorphism and reflected by it, so AOR membership
depends only on the isomorphism class of the carrier within scope
$\mathfrak S$.
\end{theorem}

\begin{proof}
The proof is the standard full-subcategory argument.  The inclusion is
identity on the underlying scoped carriers and on the morphisms between
selected objects, so it is full and faithful by construction.  The
isomorphism-invariance statement follows because each audit obligation is a
typed field-level condition: transporting a record, route, residual status,
nonclaim entry, or meta-audit discharge across an isomorphism changes its
presentation, not its existence inside the fixed scope.\footnote{Verified in Lean as \texttt{SixBirdsNeedles.AOR.PreSB.aorFullSubcat}; see \Cref{subsec:paper4-coverage-table}.}
\end{proof}
 \section{\texorpdfstring{Methods: the closure-completion functor \(\ClSB\)}{Methods: the closure-completion functor ClSB}}\label{sec:aor-methods}

\Cref{sec:aor-problem-setup} defined \(\AOR_{\mathfrak S}\) as the
full subcategory of audit-closed carriers inside
\(\PreSB_{\mathfrak S}\).  We now follow the standard reflection
pattern \citep{MacLane1998Categories,Borceux1994CategoricalAlgebra}:
construct a canonical completion of a seed carrier, show that maps
from the seed into an audit-closed carrier factor uniquely through
the completion, and package the result as a left adjoint to the
inclusion.  The construction is order-theoretic.  Over a fixed scope
\(\mathfrak S\), AOR-extending operations are ordered by subcarrier
inclusion, and a seed carrier is saturated by repeatedly adding the
minimal operations demanded by its missing audit obligations.

\begin{definition}[Closure-completion operator]\label{def:aor:cl-sb-operator}
The \emph{closure-completion operator}
$\ClSB:\PreSB_{\mathfrak S}\to\AOR_{\mathfrak S}$ assigns to a PreSB
carrier $\mathsf C$ the carrier obtained as the least fixed point of
the monotone-saturation map
\[
  \mathsf C \;\longmapsto\;
  \mathsf C \cup \{\text{minimal AOR-extending operations on }\mathsf C\}
\]
on the complete lattice of AOR-extending operations over the scope
$\mathfrak S$, together with the reflection unit
$\eta_{\mathsf C}:\mathsf C\to\ClSB(\mathsf C)$ given by the inclusion
of the seed carrier into its saturation.
\end{definition}

The lattice in the definition has as elements the candidate extensions
of the carrier's audit layer inside the fixed scope.  Its order is
subcarrier inclusion: one candidate is below another when it contains
no additional carrier or record data beyond the data of the second.
Meets and joins are computed fieldwise by intersection and generated
union of typed audit data, with the scope held fixed.  The saturation
endomap is therefore monotone: adding data before applying the rule
cannot remove any typed route, residual status, nonclaim entry, or
meta-audit discharge after applying it.

\begin{theorem}[Construction theorem for $\ClSB$]\label{thm:aor:cl-sb-construction}
The lattice of AOR-extending operations over a fixed scope
$\mathfrak S$ is a complete lattice, and the saturation map of
\Cref{def:aor:cl-sb-operator} is monotone with respect to the
subcarrier order.  By the Knaster--Tarski fixed-point theorem the
least fixed point exists, is unique, and belongs to
$\AOR_{\mathfrak S}$.  In particular $\ClSB(\mathsf C)$ is well-defined
on every PreSB carrier of scope $\mathfrak S$.
\end{theorem}

\begin{proof}
Fix the scope \(\mathfrak S\) and a seed carrier \(\mathsf C\).  Let
\(\mathcal L_{\mathfrak S}(\mathsf C)\) be the poset of AOR-extending
operation packages over \(\mathsf C\), ordered by subcarrier inclusion and
typed-field extension.  The fieldwise meet of any family keeps precisely the
data common to all packages, and the fieldwise generated join adjoins the
union of all carrier, route, residual, nonclaim, and meta-audit data while
retaining only well-typed scope-compatible entries.  These operations give a
complete lattice.  The saturation endomap \(S\) sends a package to the same
package together with the minimal operations forced by its currently missing
AOR obligations.  If \(X\le Y\), then every obligation visible in \(X\) is
also visible in \(Y\), and every operation forced over \(X\) is included
among the operations forced over \(Y\); hence \(S(X)\le S(Y)\).

Knaster--Tarski \citep{TarskiFixedPoint1955} now gives a least fixed
point \(\mathrm{lfp}(S)\).  This object contains the seed package,
because the seed is the bottom datum relative to which the lattice was
formed, and the unit \(\eta_{\mathsf C}\) is the induced inclusion.  It is AOR because
being fixed by \(S\) says that no further minimal AOR-extending operation is
missing.  Equivalently, every typed-record, route, residual-status,
nonclaim, and meta-audit obligation demanded by the package has already
been supplied.  Uniqueness is uniqueness of the least fixed point in a
poset: any other fixed point containing the seed is above \(\mathrm{lfp}(S)\).
Thus \(\ClSB(\mathsf C)\) is defined for every scoped PreSB carrier.\footnote{Mechanized in Lean as \texttt{SixBirdsNeedles.AOR.ClSB.clSBConstruction}; the Lean theorem proves the closure package, monotone saturation record, zero-defect target, and identity-field counit under the complete-lattice placeholder narrowing. The full statement is proved here. See \Cref{subsec:paper4-coverage-table}.}
\end{proof}

\begin{lemma}[Initiality of saturation steps]
Let \(\mathsf C\) be a PreSB carrier of scope \(\mathfrak S\), and let
\(\mathsf D\) be an AOR carrier of the same scope.  For every PreSB
morphism \(g:\mathsf C\to\iota\mathsf D\) and every minimal
AOR-extending operation \(\sigma\) added by the saturation map \(S\) to
a candidate package \(X\supseteq\mathsf C\), there is a unique PreSB
morphism extension
\(\widetilde g_\sigma:S(X)\to\iota\mathsf D\) such that
\(\widetilde g_\sigma|_X=g_X\).
\end{lemma}

\begin{proof}
The minimal operation \(\sigma\) adds exactly one of the audit-completion
data types used by the saturation map: typed-record data, route audits,
residual statuses, nonclaim entries, or meta-audit discharge records.  In
each case, the AOR-extending datum of \(\sigma\) is the smallest typed
completion of \(X\) that closes the obligation \(\sigma\) addresses.  The
map \(g_X\) sends the fields already present in \(X\) to a typed
configuration in \(\mathsf D\).  Because \(\mathsf D\) is AOR, the
corresponding obligation is already discharged in \(\iota\mathsf D\), so
the image of the new completion datum is forced.  This defines
\(\widetilde g_\sigma\).

Uniqueness is fieldwise.  Any extension of \(g_X\) across \(S(X)\) must
agree with \(g_X\) on \(X\), and the only new datum is the minimal
completion \(\sigma\).  Since \(\sigma\) closes the smallest obligation
not yet closed in \(X\), its image in the AOR target is determined by
the already fixed target fields.  Hence two such extensions agree on
both \(X\) and \(\sigma\), and therefore on all of \(S(X)\).
\end{proof}

\begin{theorem}[Universal property of $\ClSB$]\label{thm:aor:cl-sb-universal-property}
For every PreSB carrier $\mathsf C$, every AOR carrier $\mathsf D$ of
the same scope, and every PreSB-morphism $g:\mathsf C\to\iota\mathsf D$,
there exists a unique AOR-morphism
$\bar g:\ClSB(\mathsf C)\to\mathsf D$
satisfying $g = \iota(\bar g)\circ\eta_{\mathsf C}$.
\end{theorem}

\begin{proof}
Let \(g:\mathsf C\to\iota\mathsf D\) be a PreSB morphism and suppose
\(\mathsf D\) is already AOR in the same scope.  The carrier
\(\ClSB(\mathsf C)\) is the least fixed point of the saturation map \(S\)
starting from \(\mathsf C\).  At each successor stage of the
Knaster--Tarski construction, the initiality lemma gives a unique PreSB
morphism extension across every minimal AOR-extending operation added by
\(S\).  At a limit stage, take the fieldwise union of the compatible
extensions already constructed; compatibility holds because each successor
extension restricts to the preceding one.

Transfinite induction over the fixed-point iteration therefore gives a
unique PreSB morphism
\(\widetilde g:\ClSB(\mathsf C)\to\iota\mathsf D\) with
\(\widetilde g\circ\eta_{\mathsf C}=g\).  Since \(\ClSB(\mathsf C)\) and
\(\mathsf D\) are AOR carriers, \(\widetilde g\) is equivalently the
inclusion of a unique AOR morphism
\(\bar g:\ClSB(\mathsf C)\to\mathsf D\).  The displayed factorization
holds by construction of \(\eta_{\mathsf C}\) as the seed inclusion.

For uniqueness, let \(\bar g'\) be another AOR morphism with the same
factorization property.  Its inclusion into \(\PreSB_{\mathfrak S}\)
agrees with \(g\) on the seed.  By the initiality lemma it must agree
with the constructed extension at every successor saturation step, and
by fieldwise union it must agree at limit stages.  Hence the two maps
agree on the least fixed point \(\ClSB(\mathsf C)\), so the
factorization is unique.\footnote{Mechanized in Lean as \texttt{SixBirdsNeedles.AOR.ClSB.clSBUniversalProperty}; the Lean theorem proves the unit, lifted morphism, and fieldwise uniqueness package under the narrowed reflection-package encoding. The full statement is proved here. See \Cref{subsec:paper4-coverage-table}.}
\end{proof}

\begin{theorem}[$\ClSB\dashv\iota$ adjunction]\label{thm:aor:cl-sb-adjunction}
The closure-completion operator $\ClSB$ extends to a functor
$\PreSB_{\mathfrak S}\to\AOR_{\mathfrak S}$ which is left adjoint to
the inclusion $\iota$, with unit $\eta$ as in
\Cref{def:aor:cl-sb-operator}.  The counit is the identity on AOR
carriers: $\ClSB\circ\iota=\mathrm{id}_{\AOR_{\mathfrak S}}$.
\end{theorem}

\begin{proof}
On objects, the functor sends \(\mathsf C\) to its least saturation
\(\ClSB(\mathsf C)\).  On a morphism \(f:\mathsf C\to\mathsf C'\), apply
the universal property to the composite
\(\eta_{\mathsf C'}\circ f:\mathsf C\to \iota\ClSB(\mathsf C')\).  This
produces the unique AOR morphism
\(\ClSB(f):\ClSB(\mathsf C)\to\ClSB(\mathsf C')\) extending \(f\) through
the unit.  Identity and composition laws follow from uniqueness: both sides
of each functoriality equation are AOR morphisms whose inclusions agree
after precomposition with the seed unit, so the universal property identifies
them.

The same universal property gives the adjunction.  For
\(\mathsf C\in\PreSB_{\mathfrak S}\) and
\(\mathsf D\in\AOR_{\mathfrak S}\), precomposition with
\(\eta_{\mathsf C}\) gives the map
\[
  \mathrm{Hom}_{\AOR_{\mathfrak S}}(\ClSB(\mathsf C),\mathsf D)
  \longrightarrow
  \mathrm{Hom}_{\PreSB_{\mathfrak S}}(\mathsf C,\iota\mathsf D),
\]
and the preceding theorem gives its inverse by unique lifting.  Naturality
is again forced by uniqueness of the lifted morphism.  If the input object is
already \(\mathsf D\in\AOR_{\mathfrak S}\), the saturation step adds no new
operations, so \(\ClSB(\iota\mathsf D)=\mathsf D\) with identity field maps.
This is the identity counit.\footnote{Mechanized in Lean as \texttt{SixBirdsNeedles.AOR.ClSB.clSBAdjunction}; the Lean theorem proves the functorial extension, unit, and identity-field counit under the narrowed Hom-bijection encoding. The full statement is proved here. See \Cref{subsec:paper4-coverage-table}.}
\end{proof}

Thus \(\AOR_{\mathfrak S}\) is reflective inside
\(\PreSB_{\mathfrak S}\).  Applying \(\ClSB\) to an arbitrary scoped carrier
adds exactly the missing AOR-extending operations generated by monotone
saturation; applying it to an already audit-closed carrier returns the same
carrier through the identity-field counit.
 \section{Primary results}\label{sec:aor-primary-results}

\subsection{Structural defect and its eight strata}\label{sec:aor-primary-results-defect}

The first diagnostic layer is the structural defect \(\XiSB\): the
aggregate record of ways in which a scoped PreSB carrier has not yet
closed as an AOR carrier.  Splitting the defect into strata keeps the
failure mode visible, in the spirit of typed fault-and-failure
taxonomies in the dependable-computing literature
\citep{Avizienis2004Dependability}.  A missing observation typing, an
undeclared stage law, an unwitnessed constraint, an unannotated
rewrite, a missing holonomy transport, an unstatused residual, an
unregistered nonclaim, and an undischarged meta-audit obligation are
different audit failures, even when they all obstruct the same
conclusion that the carrier is audit-closed.  Strata are recorded as
finite multisets in the sense of \citet{Blizard1989Multiset}, so a
record may appear with multiplicity matching the carrier-field data
that creates it.

\begin{definition}[Structural defect $\XiSB$]\label{def:aor:xi-sb}
The \emph{structural defect} $\XiSB(\mathsf C)$ of a PreSB carrier
$\mathsf C$ is the disjoint sum of its eight defect strata,
\begin{align*}
  \XiSB(\mathsf C)
  \;=\;& \Sigobs(\mathsf C) \sqcup \Sigstage(\mathsf C) \sqcup \Sigconstr(\mathsf C) \sqcup \Sigrewrite(\mathsf C) \\
  & {} \sqcup \Sighol(\mathsf C) \sqcup \Sigres(\mathsf C) \sqcup \Signon(\mathsf C) \sqcup \Sigmeta(\mathsf C),
\end{align*}
with the strata defined in \Cref{def:aor:sigma-obs}--\Cref{def:aor:sigma-meta}.
The carrier is AOR exactly when every stratum is empty.
\end{definition}

The first three strata record failures in the typed interface between the
carrier and its declared audit data: observations must have declared
codomains and roles, stages must have refinement laws, and constraints must
have typed witnesses.

\begin{definition}[$\Sigobs$ stratum]\label{def:aor:sigma-obs}
The \emph{observation defect stratum} $\Sigobs(\mathsf C)$ is the
multiset of typed observation records on $\mathsf C$ whose target
codomain, role, source provenance, or evaluation context is not fully
declared.  The stratum is empty iff every declared observation in
$\mathrm{Obs}$ is a fully typed record.
\end{definition}

\begin{definition}[$\Sigstage$ stratum]\label{def:aor:sigma-stage}
The \emph{stage defect stratum} $\Sigstage(\mathsf C)$ is the multiset
of stage / refinement entries whose stage law, refinement axis, or
compatibility with the route registry is not declared.  The stratum is
empty iff every entry of $\mathrm{Stage}$ has a declared refinement
kind.
\end{definition}

\begin{definition}[$\Sigconstr$ stratum]\label{def:aor:sigma-constraint}
The \emph{constraint defect stratum} $\Sigconstr(\mathsf C)$ is the
multiset of constraint and budget entries that fail positivity, fail
the audit-pairing typing, or are paired with no declared admissible
solution.  The stratum is empty iff every entry of $\mathrm{Constr}$
admits a typed witness.
\end{definition}

The rewrite and holonomy strata capture failures of motion through the audit
system.  Rewrite entries need explicit source and target presentations;
route-dependent memory needs a typed carrier-route and declared attribution.

\begin{definition}[$\Sigrewrite$ stratum]\label{def:aor:sigma-rewrite}
The \emph{rewrite defect stratum} $\Sigrewrite(\mathsf C)$ is the
multiset of rewrite records whose source and target presentations,
admissibility tags, or blocked-rewrite annotations are not declared.
The stratum is empty iff every $\mathrm{Rewrite}$ entry is either an
allowed or an explicitly blocked rewrite.
\end{definition}

\begin{definition}[$\Sighol$ stratum]\label{def:aor:sigma-holonomy}
The \emph{holonomy defect stratum} $\Sighol(\mathsf C)$ is the multiset
of route-dependent memory records whose carrier transport, source-ledger
attribution, or nonclaim flag is not declared.  The stratum is empty
iff every $\mathrm{Hol}$ entry has a typed carrier-route.
\end{definition}

The last three strata collect unpaid residuals, implicit exclusions,
and meta-audit failures.  The residual stratum is refined in the next
subsection; for now it is enough to record that residual records must
carry both a type and a discharge status.

\begin{definition}[$\Sigres$ stratum]\label{def:aor:sigma-residual}
The \emph{residual defect stratum} $\Sigres(\mathsf C)$ is the multiset
of Schur-typed residuals on $\mathsf C$, further decomposed by the
twelve residual types of \Cref{def:aor:residual-types-twelve}.  The
stratum is empty iff every residual record carries a discharge status
from the AOR-8 catalogue of \Cref{def:aor:aor-8-statuses}.
\end{definition}

\begin{definition}[$\Signon$ stratum]\label{def:aor:sigma-nonclaim}
The \emph{nonclaim defect stratum} $\Signon(\mathsf C)$ is the
multiset of statements that the carrier needs to exclude (scope
boundaries, blocked routes, blocked rewrites, deprecated claims) but
that are not registered in $\mathrm{NC}$.  The stratum is empty iff
every implicit nonclaim of the scope has an explicit entry.
\end{definition}

\begin{definition}[$\Sigmeta$ stratum]\label{def:aor:sigma-meta}
The \emph{meta defect stratum} $\Sigmeta(\mathsf C)$ is the multiset
of meta-audit obligations (provenance trail, version pin, audit pass
log, dispute resolution) that are not discharged in $\mathrm{Meta}$.
The stratum is empty iff every meta obligation has a typed audit
record.
\end{definition}

The symbol \(\XiSB\) denotes this AOR structural defect only; it is not the
adequacy residual \(\Xi_C(D\mid L)\) of \citet{Tsiokos2026XiCompanion}.

\begin{theorem}[Eight-stratum decomposition]\label{thm:aor:eight-stratum-decomposition}
For every PreSB carrier $\mathsf C$, the structural defect
$\XiSB(\mathsf C)$ decomposes as a disjoint sum of the eight strata
\[
  \Sigobs,\ \Sigstage,\ \Sigconstr,\ \Sigrewrite,\ \Sighol,\ \Sigres,\ \Signon,\ \Sigmeta,
\]
and the decomposition is computable from the field data of
$\mathsf C$.  Each stratum is preserved by PreSB-isomorphism.
Consequently $\XiSB(\mathsf C)=0$ if and only if each of the eight
strata is independently empty.
\end{theorem}

\begin{proof}
Each defect record is created from a typed field of the carrier.  The field
that creates it also determines its stratum: an observation record belongs
to \(\Sigobs\), a stage entry to \(\Sigstage\), a constraint or budget entry
to \(\Sigconstr\), a rewrite record to \(\Sigrewrite\), a holonomy or route
memory record to \(\Sighol\), a residual record to \(\Sigres\), an implicit
exclusion to \(\Signon\), and a meta-audit obligation to \(\Sigmeta\).  These
field tags are disjoint by construction, so the structural defect is the
coproduct of the eight typed collections.  Computability is simply the
fieldwise scan of the carrier: inspect each declared record, test whether
the corresponding required annotation is present, and emit the uniquely
tagged defect record exactly when it is absent.

A PreSB isomorphism transports each typed field to the corresponding typed
field of the target carrier and transports the inverse field back.  It
therefore carries observation defects to observation defects, stage defects
to stage defects, and similarly for all eight strata.  Since the inverse
isomorphism supplies the reverse transport, emptiness of a stratum is
preserved and reflected.  Finally, a disjoint sum is empty exactly when each
summand is empty, which gives the stated criterion for \(\XiSB(\mathsf C)=0\).\footnote{Verified in Lean as \texttt{SixBirdsNeedles.AOR.DefectStrata.eightStratumDecomposition}; see \Cref{subsec:paper4-coverage-table}.}
\end{proof}

\subsection{\texorpdfstring{The twelve-type \(\Sigma_{\mathrm{residual}}\) atlas}{The twelve-type Sigma-residual atlas}}\label{sec:aor-primary-results-atlas}

The residual stratum \(\Sigres\) tracks unpaid mathematical content
rather than missing record syntax.  It is refined here along three
axes: primary type, atom refinement, and product behavior.  We call
the resulting classification the twelve-type atlas.  The displayed
list of macros names thirteen symbols because \(\Dlocal\) and
\(\Dglobal\) are presented as a paired entry; throughout the paper
they form the local--global mismatch type of the atlas.

\begin{definition}[Twelve residual types]\label{def:aor:residual-types-twelve}
The residual stratum $\Sigres$ is partitioned by twelve disjoint
\emph{residual types}: $\Ddescent$ (legal-quotient descent failure),
$\Dtransport$ (route map failure), $\Dinterface$ (bounded-interface
failure), $\Dlocal$ and $\Dglobal$ (local/global mismatch),
$\Dpresentation$ (normalization mismatch), $\Dsource$ (source-admission
failure), $\Dhorizon$ (carrier-class horizon failure), $\Dscale$
(stage-law failure), $\Dlimit$ (asymptotic failure), $\Dsym$
(symmetry/equivariance failure), $\Drole$ (role-typing failure), and
$\Dtarget$ (response-target failure).  Every $\Sigres$ entry has a
unique primary type from this list.
\end{definition}

Residual records are compared by refinement: a record refined into finer
probe or stage data is strictly smaller in the refinement order, and a
record's primary type plus its forced-secondary set fully determine the
order at the atom level.

\begin{definition}[Residual partial order]\label{def:aor:residual-partial-order}
The relation $\preceq_{\mathrm{res}}$ on $\Sigres$ is generated by:
(i) atom refinement, $a\preceq_{\mathrm{res}}b$ when $a$ is obtained
from $b$ by splitting along a finer probe family or stage; (ii) type
inclusion, $a\preceq_{\mathrm{res}}b$ when $a$ has $b$'s primary type
together with strictly more forced secondary types than $b$; and
(iii) componentwise products.  It is a partial order on isomorphism
classes of residual records.
\end{definition}

\begin{definition}[Residual atom]\label{def:aor:residual-atom}
A residual record $a\in\Sigres$ is an \emph{atom} if it is minimal
under $\preceq_{\mathrm{res}}$ among non-empty records.  Equivalently,
$a$ has a unique primary type and no proper non-empty refinement.
\end{definition}

\begin{theorem}[Atom decomposition]\label{thm:aor:atom-decomposition}
Every residual record $r\in\Sigres$ admits a finite decomposition
$r=\bigsqcup_{i} a_i$ as a disjoint sum of atoms.  The decomposition is
unique up to permutation, and is computable from the carrier field
data.
\end{theorem}

\begin{proof}
The residual registry is finite at the carrier-field level used here.  Start
with a non-empty residual record \(r\).  If \(r\) is not minimal, the
definition of \(\preceq_{\mathrm{res}}\) supplies a proper non-empty
\(c\prec_{\mathrm{res}} r\): either a split along finer probe or stage data,
an added forced-secondary component, or a componentwise product split.
Repeating this process must stop because each step moves to a strictly
smaller finite record in the carrier's declared residual data.  The terminal
pieces are precisely the atoms.

The decomposition is unique up to permutation because the refinement steps
are disjoint-sum refinements of typed records.  Two complete decompositions
therefore have the same multiset of minimal typed components: compare the
primary type and carrier-field provenance of each terminal component, match
equal components, and use minimality to rule out any further splitting on
one side but not the other.  The same field scan that detects residual
defects computes the decomposition by recursively applying the finite
refinement test.\footnote{Verified in Lean as \texttt{SixBirdsNeedles.AOR.ResidualAtlas.atomDecomposition}; see \Cref{subsec:paper4-coverage-table}.}
\end{proof}

\begin{theorem}[Residual product]\label{thm:aor:residual-product}
The product of two PreSB carriers $\mathsf C_1$ and $\mathsf C_2$ over
the same scope (componentwise product of all seventeen fields) has
residual stratum
$\Sigres(\mathsf C_1\times\mathsf C_2)
   = \Sigres(\mathsf C_1) \oplus \Sigres(\mathsf C_2)
     \oplus \Sigres^{\mathrm{cross}}(\mathsf C_1,\mathsf C_2)$,
where $\Sigres^{\mathrm{cross}}$ collects only those residuals whose
primary type is $\Dinterface$, $\Dlocal$, $\Dglobal$, or $\Dsym$ and
that arise from the product pairing.  In particular the eight other
primary types are additive across products.
\end{theorem}

\begin{proof}
The product carrier is defined componentwise across the seventeen fields.
A residual whose evidence lives entirely in the first component contributes
the corresponding element of \(\Sigres(\mathsf C_1)\); similarly, residual
evidence entirely in the second component contributes an element of
\(\Sigres(\mathsf C_2)\).  These two contributions are additive because the
field tags and carrier projections keep their provenance.

The only new residuals are those created by pairing the two components.
Such cross terms can arise only where product formation genuinely couples
the two carriers: at bounded interfaces, local/global compatibility checks,
and symmetry or equivariance constraints.  These are exactly the primary
types \(\Dinterface,\Dlocal,\Dglobal,\Dsym\).  Descent, transport,
presentation, source, horizon, scale, limit, role, and target failures are
read from a single component field and therefore do not create cross terms.
This gives the displayed decomposition and the additivity statement for the
remaining primary types.\footnote{Verified in Lean as \texttt{SixBirdsNeedles.AOR.ResidualAtlas.residualProduct}; see \Cref{subsec:paper4-coverage-table}.}
\end{proof}

\begin{theorem}[Forced-residual theorem]\label{thm:aor:forced-residual}
For each primary residual type $\Delta_*$, there is a fixed minimal set
of \emph{forced secondary types} that any non-empty record of primary
type $\Delta_*$ must also exhibit.  The complete forced-secondary list is
recorded in \Cref{tab:aor-forced-secondary-types}.  Forcing is faithful:
any discharge that closes the primary type without simultaneously closing
each forced secondary type leaves a non-empty $\Sigres$ entry.
\end{theorem}

\begin{proof}
Fix a primary residual type \(\Delta_*\) and an arbitrary non-empty record
\(r\) of primary type \(\Delta_*\).  We show that \(r\) must exhibit every
forced secondary type in the list assigned to \(\Delta_*\) by
\Cref{tab:aor-forced-secondary-types}, and that any discharge of \(r\) that
closes the primary type without closing each forced secondary leaves a
non-empty \(\Sigres\) entry.

By the residual-atom decomposition of \Cref{thm:aor:atom-decomposition},
\(r\) decomposes uniquely as a sum of atoms whose primary tags refine
\(\Delta_*\).  Each such atom carries a discharge audit obligation: it
requires the audit fields whose primary tag is \(\Delta_*\) and every forced
secondary tag listed for \(\Delta_*\).  Omitting any forced secondary tag from
the discharge leaves the field tagged with that secondary type in the
unclosed state, which by definition contributes a nonzero \(\Sigres\) entry
of that secondary type.  Thus the forced-secondary list cannot be reduced.

The forced-secondary list is derived row by row from the typed dependencies
of the audit rules.  For \(\Dsource\), the table forces
\(\Dtransport,\Dscale,\Drole,\Dtarget\), because source admission must specify
how source data are transported, scaled, role-typed, and targeted.  For
\(\Dtransport\), the table forces
\(\Dscale,\Dlimit,\Drole,\Dtarget,\Dpresentation\), because transport audits
depend on scale, limit behavior, role typing, target typing, and presentation
normalization.  For \(\Dinterface\), the table forces
\(\Dlocal,\Dglobal\), because a bounded-interface failure is visible only
relative to the local and global compatibility checks.  The remaining rows
are the same dependency-closure check, with the complete enumeration recorded
in \Cref{tab:aor-forced-secondary-types}.

Faithfulness is now immediate.  If a discharge witness for \(r\) closes
\(\Delta_*\) but omits a forced secondary type, the preceding dependency
argument leaves the corresponding secondary-tagged field unclosed.  That
field contributes a non-empty element of \(\Sigres\), so the original
residual record has not been fully discharged.\footnote{Verified in Lean as \texttt{SixBirdsNeedles.AOR.ResidualAtlas.forcedResidual}; see \Cref{subsec:paper4-coverage-table}.}
\end{proof}

\subsection{AOR-8 statuses and discharge}\label{sec:aor-primary-results-statuses}

The atlas classifies residuals by primary type.  Discharge requires a second
coordinate: a status saying how the atom is paid, blocked, cited, bounded, or
deferred to an asymptotic regime.  In the base finite-stage theory this is the
eleven-entry AOR-8 catalogue; the four asymptotic extensions are introduced
later with refinement chains.  A discharge record pairs one residual atom with
one such status and an appropriate witness payload, and a discharge witness is
the family of these records over the atoms of a residual.

\begin{definition}[AOR-8 status catalogue]\label{def:aor:aor-8-statuses}
The \emph{AOR-8 status catalogue} is the eleven-entry vocabulary of
discharge statuses for $\Sigres$ records:
\texttt{zero}, \texttt{budgeted}, \texttt{bridged}, \texttt{hidden},
\texttt{moduli}, \texttt{selected}, \texttt{probabilistic},
\texttt{out\_of\_scope}, \texttt{nonclaim}, \texttt{blocked},
\texttt{approved\_other}.
The asymptotic extensions (\Cref{def:aor:asymptotic-statuses}) add
\texttt{asymptotic\_budgeted}, \texttt{summable}, \texttt{contractive},
\texttt{eventual\_zero} for limit-stage records.
\end{definition}

The base catalogue is intentionally finite.  It separates genuinely closed
residuals from budgeted, bridged, selected, probabilistic, out-of-scope,
nonclaim, blocked, and externally approved records, so that the cascade layer
can compute downstream obligations from the pair consisting of a residual type
and a declared status.

\begin{definition}[Discharge record]\label{def:aor:discharge-record}
A \emph{discharge record} on a residual atom $a\in\Sigres$ is a triple
$(a,s,\pi)$ with $s$ an AOR-8 status (or asymptotic extension) and
$\pi$ a typed witness payload appropriate to $s$: a positivity proof
for \texttt{zero}; a finite budget bound for \texttt{budgeted}; a
bridging carrier identifier for \texttt{bridged}; a moduli space pin
for \texttt{moduli}; a selection rule for \texttt{selected}; a measure
record for \texttt{probabilistic}; a scope reference for
\texttt{out\_of\_scope} and \texttt{nonclaim}; a block reason for
\texttt{blocked}; an external-record citation for \texttt{approved\_other};
or an asymptotic-witness payload as appropriate.
\end{definition}

The payload \(\pi\) is deliberately status-dependent.  A zero record carries
different evidence from a budgeted record or an approved-other record; keeping
the payload typed prevents later cascade steps from treating all discharges as
mere labels.

\begin{definition}[Discharge witness]\label{def:aor:discharge-witness}
A \emph{discharge witness} for a residual record $r=\bigsqcup_i a_i$
is a family $(a_i,s_i,\pi_i)$ of discharge records, one per atom in
the decomposition of $r$.  The witness \emph{closes} $r$ if every
status $s_i$ is one of the AOR-8 statuses except \texttt{open}, and
every forced-secondary type of $a_i$ also has a discharge record in
the witness.
\end{definition}

The cascade-stable layer uses these definitions with a fixed
carrier-level discharge assignment.  Stability is therefore not
obtained by choosing a status to suit the atom; the assignment is part
of the audit object, and the cascade test asks whether the declared
assignment drives every residual atom to closure.
 \section{\texorpdfstring{Secondary results: discharge cascades over \(\Sigma_{\mathrm{residual}}\)}{Secondary results: discharge cascades over Sigma-residual}}\label{sec:aor-secondary-results}

\Cref{sec:aor-primary-results} classified residual failures by type
and attached statuses and witness payloads to residual atoms.  This
section records what happens after an atom is paid: the payment may
generate downstream obligations, those obligations may themselves
require statuses, and the resulting process is a cascade.  The
cascade layer is the operational part of the audit theory: it turns
the static residual atlas into a finite discharge process whenever
the declared status assignment is compatible with the cost table.

\begin{definition}[Cost record]\label{def:aor:cost-record}
For each pair $(t,s)$ with $t$ a residual type and $s$ a status, the
\emph{cost record} $\mathrm{Cost}(t,s)$ is a finite multiset of
$(t',s')$ pairs called the \emph{downstream obligations} generated by
paying a $t$-atom with status $s$.  Cost records are tabulated in the
$132$-entry cost table $\mathrm{Cost}:\mathrm{Types}\times
\mathrm{Statuses}\to\mathcal P_{\mathrm{fin}}(\mathrm{Types}\times
\mathrm{Statuses})$.
\end{definition}

The cost record is the local transition rule.  It does not say that a
residual has been globally closed; it says which new atoms and statuses must
be checked after one atom has been paid with one declared status.

\begin{definition}[Downstream operator]\label{def:aor:down-operator}
The \emph{downstream operator} $\mathrm{Down}$ acts on a residual
record $r$ with discharge witness $(a_i,s_i,\pi_i)$ by
\[
  \mathrm{Down}(r,\{(a_i,s_i,\pi_i)\})
  \;=\;
  \bigsqcup_i \mathrm{Cost}(\mathrm{type}(a_i),s_i)
\]
viewed as a new residual record of fresh atoms to be statused.
\end{definition}

Thus \(\mathrm{Down}\) is obtained by applying the local cost rule atom by atom and
taking the disjoint union of the downstream obligations.

\begin{definition}[Cost table]\label{def:aor:cost-table}
The \emph{cost table} is the function
$\mathrm{Cost}:\mathrm{Types}\times\mathrm{Statuses}\to
\mathcal P_{\mathrm{fin}}(\mathrm{Types}\times\mathrm{Statuses})$
with $12$ residual types and $11$ base AOR-8 statuses, totalling $132$
entries.  Cost-free pairs $(t,s)$ are exactly those with
$\mathrm{Cost}(t,s)=\emptyset$.
\end{definition}

The cost table is framework data, not a uniqueness theorem.  The
\(132\)-entry base table counts \(12\) residual types against the
\(11\) base AOR-8 statuses; the asymptotic limit rows belong to the
separate asymptotic-status table used in the refinement-chain
analysis and are not part of this base count.  The table is treated
under the \texttt{definition\_metadata} discipline: the body uses the
displayed representative non-trivial rows and structural defaults, and
\Cref{tab:aor-cost-table-representative} records the representative
subset used for reproducibility.  Alternative cost tables correspond
to different framework choices with different cascade behavior, not to
alternative proofs of a canonical table.

\begin{lemma}[Cost-table admissibility]\label{lem:aor:cost-table-admissibility}
The cost table of \Cref{def:aor:cost-table} satisfies two admissibility
properties:
\begin{enumerate}
  \item \emph{Strict descent.}  For every cost-table row
  \(\mathrm{Down}(a,s)=w\), the output multiset \(w\) is strictly less
  than the singleton multiset \(\{(a,s)\}\) in the Dershowitz--Manna
  multiset extension \citep{DershowitzManna1979Multiset} of the
  well-founded order on AOR-8 status alone, where statuses are ordered by
  the rank of \Cref{def:aor:aor-8-statuses} with \texttt{zero} as the unique
  minimum and \texttt{contractive}, \texttt{eventual\_zero},
  \texttt{summable}, and \texttt{asymptotic\_budgeted} ordered above
  \texttt{zero} but below the base non-terminal statuses \texttt{budgeted},
  \texttt{bridged}, \texttt{hidden}, \texttt{moduli}, \texttt{selected},
  \texttt{probabilistic}, \texttt{out\_of\_scope}, \texttt{nonclaim},
  \texttt{blocked}, \texttt{approved\_other}, with
  \texttt{summable}<\texttt{asymptotic\_budgeted} inside the asymptotic
  tier.  Concretely, every output element \((t',s')\in w\) satisfies
  \(s'<s\) in this status order; the primary type coordinate is
  unconstrained.
  \item \emph{Local diamond.}  For every record \(r\) with two distinct
  cost-table moves \(r\to r_1\) and \(r\to r_2\), there exists a common
  cost-table-reachable record \(r_*\) with \(r_1\to^* r_*\) and
  \(r_2\to^* r_*\).
\end{enumerate}
\end{lemma}

\begin{proof}
For strict descent, fix a cost-table row \(\mathrm{Down}(a,s)=w\) and verify
\(w<\{(a,s)\}\) in the Dershowitz--Manna multiset extension on AOR-8 status.
Since the multiset extension over a well-founded order is well-founded, it
suffices to show that every output element \((t',s')\in w\) satisfies the
status descent \(s'<s\); the empty case \(w=\emptyset\) is the trivial
multiset descent.  We enumerate by row type using
\Cref{tab:aor-cost-table-representative} of
\Cref{app:extended-material}.

\emph{Case A: empty-output rows.}  Any cost-table row with
\(\mathrm{Down}(a,s)=\emptyset\) is trivially strictly less than
\(\{(a,s)\}\) in Dershowitz--Manna.  These rows include ``any base type at
\texttt{zero} \(\to\emptyset\)'', ``any base type at
\texttt{approved\_other} \(\to\emptyset\)'', the asymptotic-terminal row
\((\Dlimit,\texttt{contractive})\to\emptyset\), and any other empty-output
row in the cost-table presentation.

\emph{Case B: outputs at \texttt{zero}.}  Most structural-default rows emit a
multiset whose elements are all at status \texttt{zero}.  Since
\texttt{zero} is the unique minimum of the AOR-8 status order, every such
output element \((t',\texttt{zero})\) satisfies \(\texttt{zero}<s\) for the
input status \(s\), which is non-\texttt{zero} since the input is a
non-terminal cost-table row.  This case covers the structural-default rows
for all named primary types in the cost table:
\[
\begin{gathered}
\Ddescent,\ \Dtransport,\ \Dinterface,\ \Dlocal,\ \Dglobal,\ \Dpresentation,\\
\Dsource,\ \Dhorizon,\ \Dscale,\ \Dlimit,\ \Dsym,\ \Drole,\ \Dtarget.
\end{gathered}
\]
It also covers the status-family representative rows for any non-explicit
type at \texttt{budgeted}, \texttt{bridged}, \texttt{moduli},
\texttt{selected}, \texttt{probabilistic}, \texttt{out\_of\_scope},
\texttt{nonclaim}, \texttt{blocked}, and \texttt{hidden}.  Multiset descent
follows because every output element has a status strictly less than the
input status.

\emph{Case C: non-\texttt{zero} explicit outputs.}  Three rows have outputs
at a non-\texttt{zero} status.  The transport blocked row emits
\((\Dtransport,\texttt{budgeted})\) and \((\Drole,\texttt{zero})\); the scale
hidden row emits \((\Dscale,\texttt{blocked})\); and the limit
asymptotic-budgeted row emits \((\Dlimit,\texttt{summable})\).  In each case,
every output element's status is strictly less than the input status in the
AOR-8 rank: \texttt{budgeted} and \texttt{zero} are both \(<\)
\texttt{blocked} in the transport row; \texttt{blocked} \(<\) \texttt{hidden}
in the scale row; and \texttt{summable} \(<\)
\texttt{asymptotic\_budgeted} in the limit row.  These three rows therefore
satisfy multiset descent on status.

\emph{Case D: asymptotic-terminal rows.}  The asymptotic-terminal rows, for
example the \(\Dlimit\)-contractive row with empty output and other rows of
the asymptotic-status family, are subsumed by Cases A and B above; no
additional analysis is required.

Cases A, B, C, and D exhaust the cost-table presentation, so every output
multiset satisfies the strict-descent invariant on the status order.
Multiset descent on a well-founded order is well-founded by
Dershowitz--Manna, completing the proof of clause (1).

For the local diamond, the cost-table base assigns a unique output to each
primary-type and status pair.  Two distinct moves from the same record \(r\)
must therefore arise either from two distinct atoms in \(r\), in which case
the moves act on disjoint atoms and the diamond closes by performing each
move independently on the other atom's residue, or from a single atom
carrying two assigned statuses, which is excluded by the
single-status-per-atom rule of \Cref{def:aor:discharge-record}.  In the first
case, the common record \(r_*\) is obtained by performing both moves in
either order.
\end{proof}

\begin{definition}[Cascade graph]\label{def:aor:cascade-graph}
The \emph{cascade graph} of a residual record $r$ is the bipartite
labelled multigraph $G(r)$ whose vertices are residual-record atoms
and discharge witnesses (alternating layers), and whose edges record
the downstream obligations produced by $\mathrm{Down}$.  Edges are
weighted by their AOR-8 status.
\end{definition}

The graph \(G(r)\) is a bookkeeping device for the iterated discharge
process: atom vertices ask for evidence, witness vertices supply a
status and payload, and outgoing edges list the obligations generated
by the cost table.  Bipartite labelled multigraphs are treated with
the standard graph-theory conventions of \citet{Diestel2017GraphTheory}.

\begin{theorem}[Bipartite cascade graph]\label{thm:aor:cascade-bipartite}
For every residual record $r$ and every choice of discharge witness,
the cascade graph $G(r)$ is bipartite: edges alternate between
atom-vertices and witness-vertices, and no edge connects two vertices
of the same kind.
\end{theorem}

\begin{proof}
The definition of \(G(r)\) has exactly two vertex sorts: residual atoms and
discharge witnesses.  A cascade edge is created only when an atom is paired
with a witness or when a witness emits downstream atom obligations through
\(\mathrm{Down}\).  Both constructions cross from one sort to the other.  There is
no constructor for an atom-to-atom edge or a witness-to-witness edge, so the
graph is bipartite by inspection.\footnote{\hbadness=10000\emergencystretch=2em Mechanized in Lean as
\texttt{SixBirdsNeedles.AOR.Cascades.cascadeBipartite}; the Lean
theorem proves atom/witness alternation for the encoded cascade graph
under the finite graph-record narrowing. The full statement is proved
here. See \Cref{subsec:paper4-coverage-table}.}
\end{proof}

\begin{theorem}[Cascade termination]\label{thm:aor:cascade-termination}
The cascade graph $G(r)$ is well-founded under the
status-multiset measure induced by the AOR-8 status rank and the
Dershowitz--Manna multiset extension.  In particular every cascade
terminates in finitely many discharge steps.
\end{theorem}

\begin{proof}
Assign to each active atom-status obligation its status rank and collect the
active obligations as a finite multiset of statuses.  The rank has
\texttt{zero} as the unique minimum, places terminal asymptotic statuses
below the non-terminal base statuses, and orders the explicit non-terminal
rows so that \texttt{budgeted}<\texttt{blocked}<\texttt{hidden} and
\texttt{summable}<\texttt{asymptotic\_budgeted}.

By \Cref{lem:aor:cost-table-admissibility}(1), every cost-table move maps a
singleton multiset \(\{(a,s)\}\) to a strictly smaller multiset in the
Dershowitz--Manna extension of the status order.  Multiset descent in
a well-founded order is well-founded, so no infinite descending chain of
active obligations exists.  Therefore an iterated cascade can take only
finitely many discharge steps before no active downstream obligation remains.\footnote{Mechanized in Lean as \texttt{SixBirdsNeedles.AOR.Cascades.cascadeTermination}; the Lean theorem proves the well-founded discharge relation under the status-rank measure and finite-record narrowing. The full statement is proved here. See \Cref{subsec:paper4-coverage-table}.}
\end{proof}

\begin{theorem}[Cascade confluence]\label{thm:aor:cascade-confluence}
The discharge relation on $\Sigres$ is locally confluent: if a residual
atom admits two distinct one-step discharges into records $r_1$ and
$r_2$, then there exists a record $r_3$ reached by further discharges
from both $r_1$ and $r_2$.  Together with termination
(\Cref{thm:aor:cascade-termination}), Newman's lemma
\citep{Newman1942Combinatorial} yields global confluence
\citep{Huet1980ConfluentReductions}: every cascade has a unique normal
form.
\end{theorem}

\begin{proof}
By \Cref{lem:aor:cost-table-admissibility}(2), the local diamond holds for
every pair of distinct cost-table moves.  By
\Cref{lem:aor:cost-table-admissibility}(1) and
\Cref{thm:aor:cascade-termination}, the discharge relation is well-founded:
each cost-table move gives strict multiset descent in the
Dershowitz--Manna extension of the status order.

Newman's lemma \citep{Newman1942Combinatorial} states that local confluence
together with well-founded termination implies global confluence and unique
normal forms.  Applying it to the cost-table rewrite relation gives the
global confluence and unique-normal-form conclusion for discharge cascades.\footnote{Mechanized in Lean as
\texttt{SixBirdsNeedles.AOR.Cascades.cascadeConfluence} as a partial
mechanization: local confluence and the well-founded discharge relation are
recorded in Lean; the Newman implication and the resulting
unique-normal-form conclusion are supplied here as a classical input. See
\Cref{subsec:paper4-coverage-table}.}
\end{proof}

\begin{theorem}[Cross-cascade composition]\label{thm:aor:cross-cascade-composition}
For two residual records $r_1$ and $r_2$ on the same carrier, their
disjoint-union cascade $G(r_1\sqcup r_2)$ equals the disjoint union of
cascades $G(r_1)\sqcup G(r_2)$ together with the cross-edges generated
by interaction atoms of primary type $\Dinterface$, $\Dlocal$,
$\Dglobal$, or $\Dsym$ produced by the residual-product theorem
(\Cref{thm:aor:residual-product}).
\end{theorem}

\begin{proof}
The cascade construction is additive over disjoint unions except where the
residual-product theorem creates a cross residual.  Thus atom and witness
vertices inherited from \(r_1\) generate the cascade \(G(r_1)\), and those
inherited from \(r_2\) generate \(G(r_2)\).  The only additional edges arise
from product-pairing atoms.  By the residual-product theorem these cross
atoms have primary type \(\Dinterface,\Dlocal,\Dglobal\), or \(\Dsym\), and
the corresponding downstream edges are exactly the displayed cross-edges.\footnote{Verified in Lean as \texttt{SixBirdsNeedles.AOR.Cascades.crossCascadeComposition}; see \Cref{subsec:paper4-coverage-table}.}
\end{proof}

\begin{definition}[Cascade-stable AOR]\label{def:aor:cascade-stable-aor}
A PreSB carrier $\mathsf C\in\AOR_{\mathfrak S}$ is
\emph{cascade-stable} if every residual record in $\Sigres(\mathsf C)$
admits a discharge witness whose cascade reaches the empty record in
finitely many steps using only the base AOR-8 statuses
(\Cref{def:aor:aor-8-statuses}).  The class of cascade-stable AOR
carriers is denoted $\CascadeStableAOR_{\mathfrak S}$.
\end{definition}

The predicate is evaluated on an AOR carrier together with its fixed
residual-discharge assignment.  Cascade stability is therefore not
proved by choosing a fresh \texttt{zero} status when an atom appears;
it is verified by checking that the declared assignment drives the
cost-table walk to the empty record.

\begin{theorem}[Cascade-stable characterisation]\label{thm:aor:cascade-stable-characterisation}
$\mathsf C\in\CascadeStableAOR_{\mathfrak S}$ if and only if every
$\Sigres$-atom of $\mathsf C$ has a declared discharge status whose
cost-table walk reaches a cost-free terminal status in finitely many
moves.  Equivalently, the $\omega$-iteration of $\mathrm{Down}$ from
the initial atoms of $\mathsf C$ produces an eventually empty sequence
of records.
\end{theorem}

\begin{proof}
Assume first that the carrier with its fixed discharge assignment is
cascade-stable.  Apply the definition to a residual record consisting of a
single atom.  Its declared status must produce a finite cascade to the empty
record.  Reading the successive atom-status obligations along that finite
cascade gives a finite path in the cost table from the assigned status to a
cost-free terminal status, equivalently a finite \(\omega\)-iteration of
\(\mathrm{Down}\) ending at the empty record.

Conversely, suppose every residual atom has an assigned status whose
cost-table walk reaches the empty record, and suppose the \(\omega\)-iteration
from the initial atoms is eventually empty.  Decompose an arbitrary residual
record into atoms by \Cref{thm:aor:atom-decomposition}.  For each atom use
the assumed finite cost-table walk as its discharge witness.  Taking the
disjoint union of these finite witnesses gives a witness for the whole
residual record.  Termination ensures that the combined cascade is still
finite, and the eventual-emptiness assumption says that no downstream
obligation remains.  Hence the carrier is cascade-stable.\footnote{Mechanized in Lean as \texttt{SixBirdsNeedles.AOR.Cascades.cascadeStableCharacterisation}; the Lean theorem proves the assignment-based cost-table reachability characterization under the carrier-with-discharge wrapper narrowing. The full statement is proved here. See \Cref{subsec:paper4-coverage-table}.}
\end{proof}
 \section{Analysis: refinement chains, colimits, and the hierarchy}\label{sec:aor-analysis}

Finite cascades close residuals stage by stage.  Refinement chains
introduce a second phenomenon: a carrier may be closed at every
finite stage while a residual appears only at the directed colimit.
This section extends the finite cascade story by recording refinement
morphisms, directed colimits in \(\RefPreSB\), and four asymptotic
statuses for limit-stage residuals.  The asymptotic statuses are not
additional base AOR-8 statuses; they form a separate limit-stage
layer used only when colimits enter the audit problem.

\begin{definition}[Refinement morphism]\label{def:aor:refinement-morphism}
A \emph{refinement morphism} $\rho:\mathsf C\to\mathsf C'$ between two
PreSB carriers is a PreSB morphism that additionally records, for each
field, whether it is an A1--A6 refinement step
(\Cref{def:aor:refinement-kinds-six}).  Refinement morphisms compose
and have identities.  The category $\RefPreSB_{\mathfrak S}$ has PreSB
carriers as objects and refinement morphisms as arrows.
\end{definition}

A refinement morphism is therefore a typed PreSB morphism together
with an audit annotation that records the kind of refinement
performed in each field.

\begin{definition}[Six refinement kinds]\label{def:aor:refinement-kinds-six}
A refinement morphism's per-field annotations belong to one of six
disjoint kinds:
\textbf{A1} scope expansion (enlarging $\mathfrak S$ within a fixed
scope chain);
\textbf{A2} record-field expansion (adding admissible records to one
of the registries);
\textbf{A3} stage / refinement-axis growth (extending $\mathrm{Stage}$
by a new stage);
\textbf{A4} constraint tightening (replacing a constraint by a
strictly stronger one);
\textbf{A5} morphism strengthening (replacing a morphism by a
refinement of itself);
\textbf{A6} carrier (probe / target) refinement (replacing $(H,q,Q)$
or a probe/target family by a finer version, with a declared
projection map).
\end{definition}

The six kinds separate changes of scope, changes of records, changes
of staging, constraint strengthening, morphism strengthening, and
carrier refinement proper.  The separation lets a refinement chain
record whether a new limit residual is caused by a new stage, a
stronger constraint, or a finer carrier presentation.

\begin{definition}[$\RefPreSB$ category]\label{def:aor:refpresb-category}
The category $\RefPreSB_{\mathfrak S}$ is the full category of PreSB
carriers of scope $\mathfrak S$ with refinement morphisms.  A
\emph{refinement chain} is a directed diagram in
$\RefPreSB_{\mathfrak S}$ indexed by a directed poset.
\end{definition}

\begin{theorem}[$\RefPreSB$ is directed-cocomplete]\label{thm:aor:refpresb-cocomplete}
The category $\RefPreSB_{\mathfrak S}$ has all directed colimits.
Concretely, for a directed diagram
$(\mathsf C_j,\rho_{jj'})$ with projections
$p_{jj'}:\mathsf C_{j'}\to\mathsf C_j$ along reverse arrows, the
colimit object $\mathsf C_\infty=\mathrm{colim}_j\mathsf C_j$ has
record/scope data as the directed union of the corresponding $\mathsf
C_j$ data, and carrier data as the inverse limit of compatible threads
under the projections $p_{jj'}$.  The construction is functorial and
preserves PreSB-typing.
\end{theorem}

\begin{proof}
The directed-colimit construction follows the standard pattern in
accessible and locally presentable categories
\citep{AdamekRosicky1994LocallyPresentable,MacLane1998Categories}.  For
the record-like fields, directedness lets one take the union of all
records that appear at some finite stage.  Compatibility of the
refinement maps ensures that a record admitted at one stage remains
interpretable at later stages.  For the carrier-like fields, the
reverse projections \(p_{jj'}\) define the compatible-thread condition,
and the carrier part of the colimit is the inverse limit of those
threads.  The structure maps are the evident inclusions on record data
and projections on compatible carrier threads.  Functoriality follows
by applying the same union and compatible-thread constructions to
morphisms of directed diagrams, and PreSB-typing is preserved because
every finite-stage typing condition is stable under the refinement
morphisms.\footnote{Mechanized in Lean as \texttt{SixBirdsNeedles.AOR.Refinement.refpresbCocomplete}; the Lean theorem proves existence of a representative colimit witness under the base-object cocone narrowing. The full statement is proved here. See \Cref{subsec:paper4-coverage-table}.}
\end{proof}

\begin{definition}[Asymptotic statuses]\label{def:aor:asymptotic-statuses}
For a residual atom $a$ of a colimit object
$\mathsf C_\infty=\mathrm{colim}_j\mathsf C_j$, the
\emph{asymptotic statuses} extending AOR-8 are:
\texttt{asymptotic\_budgeted} (there exists a uniform finite budget
$B$ and a tail index $J_0$ with the per-stage cost bounded by $B$ for
all $j\ge J_0$);
\texttt{summable} (the per-stage costs are summable across the
directed chain);
\texttt{contractive} (the per-stage cost ratios tend to a constant
strictly less than $1$);
\texttt{eventual\_zero} (the per-stage cost is zero for all $j\ge J_0$
for some tail index $J_0$).
\end{definition}

These four statuses are limit-stage statuses; they do not alter the
eleven-entry AOR-8 catalogue used for finite cascade discharge, but
specify how a residual of primary type \(\Dlimit\) is paid at the
colimit.

\begin{theorem}[Asymptotic soundness]\label{thm:aor:asymptotic-soundness}
The four asymptotic statuses of
\Cref{def:aor:asymptotic-statuses} satisfy the implication chain
\[
  \texttt{eventual\_zero} \Rightarrow \texttt{contractive}
  \Rightarrow \texttt{summable} \Rightarrow \texttt{asymptotic\_budgeted}.
\]
Each asymptotic status closes the colimit residual at level $\Dlimit$
without producing a fresh atom in the limit.  Moreover, if every atom
of $\mathsf C_\infty$ has an asymptotic status from this list and every
finite stage $\mathsf C_j$ is in $\CascadeStableAOR_{\mathfrak S}$,
then $\mathsf C_\infty\in\AOR_{\mathfrak S}$.
\end{theorem}

\begin{proof}
The implication chain is the standard comparison among tail conditions
of non-negative series \citep{Rudin1976Principles}.  If the per-stage
cost is eventually zero, then on the same tail the cost ratios are
eventually zero and hence bounded by a constant strictly less than
one.  A contractive tail gives a summable geometric majorant after
enlarging the initial finite segment if necessary.  A summable
non-negative tail has bounded partial sums, and therefore supplies a
uniform finite budget.  This proves
\[
  \texttt{eventual\_zero} \Rightarrow \texttt{contractive}
  \Rightarrow \texttt{summable} \Rightarrow \texttt{asymptotic\_budgeted}.
\]

For closure at level \(\Dlimit\), each of the four statuses gives a typed
tail witness that prevents a fresh limit atom from being generated: eventual
zero leaves no tail cost, contractive and summable tails have finite total
downstream obligation, and asymptotic boundedness gives a uniform budget
record for the limit-stage atom.  If every finite stage is already
cascade-stable, all finite-stage residuals discharge by the previous section.
If every atom that survives to the colimit carries one of the four
asymptotic statuses, then the only possible new \(\Dlimit\) residuals are
also paid.  Thus the colimit carrier has no unpaid residual and remains in
\(\AOR_{\mathfrak S}\).\footnote{Mechanized in Lean as \texttt{SixBirdsNeedles.AOR.Refinement.asymptoticSoundness}; the Lean theorem proves the asymptotic-status implication chain and limit-residual closure predicates under the representative colimit narrowing. The full statement is proved here. See \Cref{subsec:paper4-coverage-table}.}
\end{proof}

\begin{definition}[Refinement-stable AOR]\label{def:aor:refinement-stable-aor}
A PreSB carrier $\mathsf C\in\AOR_{\mathfrak S}$ is
\emph{refinement-stable} if AOR-membership is preserved by every
refinement morphism into $\mathsf C$ and by every directed colimit of
refinement chains terminating at $\mathsf C$.  The class of
refinement-stable AOR carriers is denoted
$\RefStableAOR_{\mathfrak S}$.
\end{definition}

As with cascade stability, the formal predicate is evaluated on a carrier
paired with a declared discharge assignment.  Refinement stability adds the
requirement that limit-stage residuals are paid by the asymptotic layer.

\begin{theorem}[Refinement-stable characterisation]\label{thm:aor:refinement-stable-characterisation}
$\mathsf C\in\RefStableAOR_{\mathfrak S}$ if and only if (i)
$\mathsf C\in\CascadeStableAOR_{\mathfrak S}$ and (ii) every residual
atom of $\mathsf C$ that survives a colimit step is paid by an
asymptotic status from \Cref{def:aor:asymptotic-statuses}.
Equivalently, every refinement chain terminating at $\mathsf C$ has a
discharge witness in which limit-stage atoms carry asymptotic
statuses.
\end{theorem}

\begin{proof}
Suppose first that \(\mathsf C\) is refinement-stable.  Applying the
definition to the constant refinement chain shows that the finite-stage
carrier must be cascade-stable: otherwise the constant chain would already
contain an unpaid finite residual.  Applying the definition to an arbitrary
directed chain terminating at \(\mathsf C\) shows that any residual atom that
survives only at the colimit must be paid there; by the preceding definition,
the available colimit payments are exactly the four asymptotic statuses.

Conversely, assume that \(\mathsf C\) is cascade-stable and that every
colimit-surviving atom is paid by an asymptotic status.  Along any refinement
chain terminating at \(\mathsf C\), finite-stage atoms are discharged by
cascade stability.  The only additional atoms that can appear at the colimit
are the surviving limit-stage atoms, and by hypothesis they carry one of the
asymptotic statuses.  By \Cref{thm:aor:asymptotic-soundness}, those statuses
close the \(\Dlimit\) residual without producing a fresh atom in the limit.
Hence AOR-membership is preserved under the relevant refinement morphisms and
directed colimits.  This is precisely refinement stability.\footnote{Mechanized in Lean as \texttt{SixBirdsNeedles.AOR.Refinement.refinementStableCharacterisation}; the Lean theorem proves the cascade-stable plus colimit-payment characterization under the carrier-with-discharge and asymptotic-status narrowing. The full statement is proved here. See \Cref{subsec:paper4-coverage-table}.}
\end{proof}

\begin{theorem}[Hierarchy witnesses]\label{thm:aor:hierarchy}
There exist two named carriers
\(\mathsf C_{\mathrm{witness},1}\) and
\(\mathsf C_{\mathrm{witness},2}\), together with typed obstruction
markers, that witness the failure of the stronger stability conditions:
\(\mathsf C_{\mathrm{witness},1}\) lies in
\(\CascadeStableAOR_{\mathfrak S}\) and carries a typed obstruction marker
certifying that it does not satisfy the refinement-stability predicate;
\(\mathsf C_{\mathrm{witness},2}\) lies in \(\AOR_{\mathfrak S}\) and
carries a typed obstruction marker certifying that it does not satisfy the
cascade-stability predicate.
\end{theorem}

\begin{proof}
The first carrier \(\mathsf C_{\mathrm{witness},1}\) is constructed as in
\Cref{tab:aor-counterexample-carriers}: it has a finite cascade that closes
with base AOR-8 statuses, hence it lies in
\(\CascadeStableAOR_{\mathfrak S}\), and it carries a typed
\path{LimitAsymptoticNonpayment} marker that records a limit-stage residual
whose discharge is not closed under the refinement chain demanded by
refinement stability.  The marker is the typed witness that the carrier
fails the refinement-stability predicate; the predicate's full
quantifier-level negation is a separate framework-level verification and is
not produced here.

The second carrier \(\mathsf C_{\mathrm{witness},2}\) is constructed
similarly: it has AOR field data, hence lies in \(\AOR_{\mathfrak S}\), and
it carries a typed \path{BlockedCascadeNontermination} marker that records a
residual atom whose discharge cascade does not close in finite steps under
any AOR-8 status.  This marker is the typed witness for failure of cascade
stability.

Both carriers and their typed markers are recorded in
\Cref{tab:aor-counterexample-carriers}.\footnote{Mechanized in Lean as
\texttt{SixBirdsNeedles.AOR.Refinement.hierarchy} as a partial mechanization:
the two witness carriers and their typed obstruction markers are recorded in
Lean; full predicate-level non-membership verification remains framework
future work and is stated as scope in the abstract. See
\Cref{subsec:paper4-coverage-table}.}
\end{proof}
 

\section{Status, scope, and nonclaim boundary}\label{sec:aor-discussion}

The preceding sections develop AOR as a meta-theory of audit-closed
carriers.  Its formal objects are scope records, PreSB carriers,
closure completion, defect strata, residual atoms, discharge
cascades, and refinement chains; they are not a direct model of any
particular analytic problem.  The Navier--Stokes smooth-data carrier
appears as an AOR instance in \citet{Tsiokos2026NavierStokes}, but
that paper is a consumer of the meta-theory developed here, not a
theorem-level dependency of it.

\begin{nonclaim}[No strong-realism reading of AOR]\label{nonclaim:aor:no-strong-realism}
Membership in $\AOR_{\mathfrak S}$ is not a metaphysical or
philosophical realism claim about the carrier's underlying domain.  It
is a record-level audit predicate: every declared record is typed,
every cited route is audited, every residual is statused, and every
nonclaim is registered, within the declared scope $\mathfrak S$.  Any
extra-scope interpretation requires its own audit and is outside the
present framework.
\end{nonclaim}

The nonclaim fixes the intended reading of the full subcategory
$\AOR_{\mathfrak S}$.  A carrier belongs to $\AOR_{\mathfrak S}$ when
its declared audit obligations close inside its declared scope.  This
is a structural predicate on a typed record, not a realism statement
about the mathematical or physical domain from which that record was
drawn.

Several further scope decisions follow from the same discipline.  First,
the present paper does not develop downstream applications beyond the
audit-closed-carrier framework itself.  The current cited consumer is
the NS-as-AOR-instance construction in
\citet{Tsiokos2026NavierStokes}; any further application requires its
own scope record, residual atlas, discharge assignment, and nonclaim
ledger.

Second, the cost table used in \Cref{sec:aor-secondary-results} is
framework data rather than a uniqueness theorem.  The 132-entry base
table records how residual type and base AOR-8 status generate
downstream obligations; the asymptotic table for refinement limits is a
separate layer.  Alternative cost tables would define related but
different cascade systems, so the table should be read as part of the
chosen AOR discipline, not as a canonical classification theorem.

Third, the hierarchy statement \Cref{thm:aor:hierarchy} is a witness
theorem rather than a complete classification theorem.  It exhibits
named carriers and typed obstruction markers at the two stability
boundaries in
$\RefStableAOR \subseteq \CascadeStableAOR \subseteq \AOR$, while full
predicate-level negated-membership proofs at every level remain
framework future work.

Fourth, the cascade-confluence statement
\Cref{thm:aor:cascade-confluence} uses a classical confluence input.
The local diamond and the well-founded termination measure are proved
inside the cascade calculus; Newman's lemma then supplies global
confluence and unique normal form for the induced discharge relation.

Several adjacent lines of work in computer science treat audit data
or proof obligations as first-class.  Hoare logic
\citep{Hoare1969AxiomaticBasis} and the refinement
calculus \citep{BackVonWright1998Refinement} formalize verification
conditions and stepwise refinement, while audited-computation calculi
\citep{BaveraBonelli2018AuditedComputation,RicciottiCheney2018ExplicitAuditing}
and provenance-as-dependency semantics
\citep{CheneyAhmedAcar2011Provenance} carry typed audit histories
through computations.  AOR shares with these the discipline of
treating audit content as part of the object rather than an external
side condition, but its meta-theoretic question is different: rather
than typing executions, it asks when a typed carrier with declared
audit data admits a categorical completion into the audit-closed
full subcategory.

Finally, the cross-paper boundary is motivational and contextual
only.  The formed-layer bookkeeping of
\citet{Tsiokos2026NeedleKiller} and the adequacy-residual calculus of
\citet{Tsiokos2026XiCompanion} motivate the audit vocabulary, but the
AOR construction does not import their internal theorems as
hypotheses; \citet{Tsiokos2026NavierStokes} uses the AOR meta-theory
as a consumer rather than as an input.  The symbol $\XiSB$ continues
to denote the structural defect of a PreSB carrier, not the adequacy
residual $\Xi_C(D\mid L)$ of \citet{Tsiokos2026XiCompanion}.
 

\section{Conclusion}\label{sec:aor-conclusion}

The central result of this paper is a closure theorem for typed
audit carriers.  Every carrier in $\PreSB_{\mathfrak S}$ admits a
completion into the audit-closed full subcategory through the
reflection $\ClSB\dashv\iota$.  The construction does not merely
adjoin missing data; it makes audit closure a structural operation,
so that records, routes, residuals, statuses, nonclaims, and
meta-audits remain visible as typed fields rather than being
collapsed into an external well-formedness assertion.  In that
setting, the chain
$\RefStableAOR\subseteq\CascadeStableAOR\subseteq\AOR$ organizes
ordinary record-level audit closure, finite cascade closure, and
closure stable under refinement limits.

The framework comprises four layers.  The first is the
closure-completion reflection and its universal property.  The second
is the eight-stratum decomposition of the structural defect $\XiSB$,
which turns failure of audit closure into typed diagnostic data.  The
third is the residual layer: a twelve-type atlas, the status
catalogue, cost records, downstream obligations, and cascade graphs
tracing how residual atoms are paid.  The fourth is the refinement
layer, where directed chains and asymptotic statuses distinguish
finite discharge from discharge that only closes at a limit.
Together they form a self-contained meta-theory of audit-closed
carriers.

The scope boundary is as important as the positive construction.
\Cref{sec:aor-discussion} records that membership in
$\AOR_{\mathfrak S}$ is not a strong-realism claim about the carrier's
domain.  The paper also does not host a PDE-specific worked example;
the NS-as-AOR-instance presentation belongs to
\citet{Tsiokos2026NavierStokes}.  The cascade-confluence theorem
proves the local diamond and well-founded termination, then applies
Newman's lemma to obtain global confluence and unique normal form.  The
hierarchy theorem exhibits named witness carriers and typed
obstruction markers, while full predicate-level non-membership proofs
at every level remain framework future work.

The meta-theory developed here is intended for use by downstream
audited applications.  The first such application is the
NS-as-AOR-instance construction of
\citet{Tsiokos2026NavierStokes}; further applications require their
own scope records, residual atlases, cost tables, and discharge
assignments.  Several natural strengthenings remain open: full
predicate-level non-membership verification for the hierarchy witnesses,
additional cost-table variants with their induced cascade systems, and
data-bearing R-record extensions in the Navier--Stokes formal layer.  These are
future refinements of the audit discipline, not assumptions needed
for the closure-completion framework itself.
 
\appendix
\section{Definitions and gate predicates}\label{app:definitions}

This appendix collects the recurring AOR vocabulary used in the body.
It does not restate the formal definitions; each entry gives a short
reminder and points back to the body location where the term is
introduced.

\subsection{Definitions}\label{subsec:definitions-list}

\begin{description}[style=nextline]
\item[PreSB carrier.]
A typed carrier together with its audit data: quotient, audit form,
record registries, stages, observations, constraints, rewrites,
holonomy, residuals, statuses, nonclaims, and meta-audits.  See
\Cref{def:aor:presb-carrier}.

\item[Scope record.]
The declaration of the carrier class, source ledger, interface family,
and route registry inside which audit closure is evaluated.  See
\Cref{def:aor:scope-record}.

\item[PreSB morphism.]
A scope-preserving field-tuple morphism between PreSB carriers,
respecting the quotient, audit, route, and registry structure.  See
\Cref{def:aor:morphism-presb}.

\item[AOR subcategory.]
The full subcategory of PreSB carriers whose declared records, routes,
residuals, nonclaims, and meta-audits close inside the declared scope;
equivalently, the structural defect is zero.  See
\Cref{def:aor:aor-subcategory}.

\item[Closure completion functor $\ClSB$.]
The closure operation sending a PreSB carrier to its audit-closed
completion and serving as the left adjoint to the inclusion
$\iota:\AOR_{\mathfrak S}\hookrightarrow\PreSB_{\mathfrak S}$.  See
\Cref{def:aor:cl-sb-operator}.

\item[Structural defect $\XiSB$.]
The aggregate defect object collecting all audit-obligation failures of
a PreSB carrier as a disjoint sum of typed strata.  See
\Cref{def:aor:xi-sb}.

\item[Observation stratum $\Sigobs$.]
The defect stratum for observation records whose target space,
boundedness, or declared measurement role is not typed.  See
\Cref{def:aor:sigma-obs}.

\item[Stage stratum $\Sigstage$.]
The defect stratum for stage data whose refinement axis, monotonicity
law, or transition record is missing.  See
\Cref{def:aor:sigma-stage}.

\item[Constraint stratum $\Sigconstr$.]
The defect stratum for constraints stated without a typed predicate,
witness domain, or compatibility proof.  See
\Cref{def:aor:sigma-constraint}.

\item[Rewrite stratum $\Sigrewrite$.]
The defect stratum for rewrite or normalization steps that lack a
declared source, target, or preservation witness.  See
\Cref{def:aor:sigma-rewrite}.

\item[Holonomy stratum $\Sighol$.]
The defect stratum for route-comparison or loop data whose holonomy
record is absent or unaudited.  See \Cref{def:aor:sigma-holonomy}.

\item[Residual stratum $\Sigres$.]
The defect stratum for residual atoms that still require status,
discharge, or cascade accounting.  See
\Cref{def:aor:sigma-residual}.

\item[Nonclaim stratum $\Signon$.]
The defect stratum for exclusions or limitations that are being used
implicitly but have not been registered as nonclaims.  See
\Cref{def:aor:sigma-nonclaim}.

\item[Meta-audit stratum $\Sigmeta$.]
The defect stratum for higher-level audit obligations, such as
manifest or provenance obligations, that have not themselves been
discharged.  See \Cref{def:aor:sigma-meta}.

\item[Residual-type atlas.]
The named partition of $\Sigres$ into primary residual types such as
$\Ddescent$, $\Dtransport$, $\Dinterface$, $\Dsource$, $\Dscale$,
$\Dlimit$, $\Drole$, and $\Dtarget$.  See
\Cref{def:aor:residual-types-twelve}.

\item[Residual partial order.]
The refinement order on residual atoms: splitting an atom into finer
probe, stage, or obstruction data makes it strictly smaller in the residual
order, so an atom is a minimal non-empty record.  See
\Cref{def:aor:residual-partial-order}.

\item[Residual atom.]
A typed residual obligation consisting of a primary residual type,
forced secondary types, and provenance data tying it to the carrier's
record fields.  See \Cref{def:aor:residual-atom}.

\item[AOR-8 status catalogue.]
The eleven base statuses used to pay residual atoms, together with the
four asymptotic extensions used for refinement limits.  See
\Cref{def:aor:aor-8-statuses,def:aor:asymptotic-statuses}.

\item[Discharge record and discharge witness.]
A discharge record pairs a residual atom with a status and a witness;
a discharge witness is the family of such records used for a carrier or
cascade step.  See
\Cref{def:aor:discharge-record,def:aor:discharge-witness}.

\item[Cost record, downstream operator, and cost table.]
The cost record $\mathrm{Cost}(t,s)$ lists downstream obligations
generated by paying a residual type $t$ with status $s$; the downstream
operator forms the next residual record, and the cost table packages
the framework's rows.  The cost table is framework data under the
documented \emph{definition metadata} discipline.  See
\Cref{def:aor:cost-record,def:aor:down-operator,def:aor:cost-table}.

\item[Cascade graph.]
The bipartite graph whose layers alternate residual atoms and
discharge witnesses, with edges recording downstream obligations
created by the cost table.  See \Cref{def:aor:cascade-graph}.

\item[Cascade-stable AOR.]
The predicate that a carrier, paired with a fixed discharge assignment,
has all residual atoms paid by finite cost-table cascades.  See
\Cref{def:aor:cascade-stable-aor}.

\item[Refinement morphism, six kinds, and $\RefPreSB$.]
A refinement morphism is a PreSB morphism annotated by the A1--A6
refinement kind of each field change; $\RefPreSB$ is the resulting
category of PreSB carriers and refinement arrows.  See
\Cref{def:aor:refinement-morphism,def:aor:refinement-kinds-six,def:aor:refpresb-category}.

\item[Refinement-stable AOR.]
The refinement-stable subclass of cascade-stable AOR carriers whose
limit-stage residuals are paid by asymptotic statuses along directed
refinement chains.  See \Cref{def:aor:refinement-stable-aor}.
\end{description}

\subsection{Gate predicates}\label{subsec:gate-predicates-list}

The following five carrier-level predicates are the gate conditions in
the AOR subcategory definition.  Their conjunction is the body-level
form of the assertion that the eight-stratum structural defect $\XiSB$
vanishes; compare \Cref{def:aor:aor-subcategory} with
\Cref{thm:aor:eight-stratum-decomposition}.

\begin{description}[style=nextline]
\item[\texttt{declaredRecordsTyped}.]
Every declared observation, source, interface, route, constraint, and
related record has a specified type and carrier-level role.

\item[\texttt{citedRoutesAudited}.]
Every route or transport map cited by the carrier appears in the route
registry with the audit data needed to compare its source and target
records.

\item[\texttt{residualsStatused}.]
Every residual atom in the residual registry has an assigned AOR-8
status and therefore enters the discharge calculus.

\item[\texttt{nonclaimsRegistered}.]
Every exclusion or limitation used by the carrier is explicitly
registered as a nonclaim rather than left implicit.

\item[\texttt{metaAuditDischarged}.]
Every meta-level audit obligation, including provenance and manifest
obligations, has a typed record witnessing its discharge.
\end{description}
 \section{Methods, artifacts, and reproducibility contract}\label{app:methods}

This appendix expands the method sketches used in the body.  Its role
is reproducibility and audit discipline: it explains what data are held
fixed, what order-theoretic or graph-theoretic method is being used, and
which framework choices downstream users must preserve.  The canonical
definitions and theorem statements remain in the body.

\subsection{\texorpdfstring{Knaster--Tarski monotone saturation for
\(\ClSB\)}{Knaster--Tarski monotone saturation for ClSB}}
\label{subsec:aor-methods-clsb-saturation}

\Cref{sec:aor-methods} constructs \(\ClSB\) by applying a monotone
saturation operation to a scoped PreSB carrier.  The reproducibility
object is the complete lattice of AOR-extending operations over the
fixed scope \(\mathfrak S\).  An element of this lattice is a candidate
audit extension of the seed carrier: it adds typed record data, route
audits, residual statuses, nonclaim entries, and meta-audit discharge
records, but it does not change the declared scope.  The order is
fieldwise extension, so meets and joins are computed by fieldwise
intersection and generated union subject to scope compatibility.

The saturation map adds the minimal operations forced by the currently
visible audit gaps.  Monotonicity is the key reproducibility check: if
one candidate extension contains another, then every obligation visible
in the smaller candidate is visible in the larger one, and saturation
cannot remove already-typed audit data.  Knaster--Tarski
\citep{TarskiFixedPoint1955} then supplies a least fixed point of the
saturation map, obtained by the constructive ordinal iteration of
\citet{CousotCousotConstructive1979}.  That fixed point is the
closure-completed carrier, and the inclusion of the seed into the
fixed point is the reflection unit.

For carriers already in \(\AOR_{\mathfrak S}\), no additional
AOR-extending operation is forced.  The counit of the reflection is
therefore the identity-field counit: applying \(\ClSB\) to an already
audit-closed carrier returns the same carrier up to the fieldwise
identity package.

\subsection{Well-founded measure for cascade termination}
\label{subsec:aor-methods-cascade-measure}

\Cref{thm:aor:cascade-termination} uses a well-founded measure on
discharge states.  The state consists of the multiset of residual atoms
currently awaiting discharge, each tagged with its primary type and
AOR-8 status.  The measure is the multiset of statuses obtained from
that state by forgetting the primary-type coordinate; the type
coordinate plays no role in the termination argument.

The status order is the AOR-8 status rank of
\Cref{def:aor:aor-8-statuses} with \texttt{zero} as the unique minimum
and the asymptotic-terminal statuses placed above \texttt{zero} but
below the base non-terminal statuses, with \(\texttt{summable} <
\texttt{asymptotic\_budgeted}\) inside the asymptotic tier.  Cascade
moves are then well-founded by Dershowitz--Manna multiset descent
\citep{DershowitzManna1979Multiset} on this status order: by
\Cref{lem:aor:cost-table-admissibility}(1), every cost-table move
replaces an atom of status \(s\) by a multiset of atoms whose statuses
are each strictly less than \(s\), so the cascade's status multiset
strictly decreases at every step.  This is the method behind the
finite-termination statement of \Cref{thm:aor:cascade-termination},
and it supplies the well-founded discharge relation used as the
termination input for the confluence argument in
\Cref{thm:aor:cascade-confluence}.

\subsection{\texorpdfstring{Representative-object cocompleteness for
\(\RefPreSB\)}{Representative-object cocompleteness for RefPreSB}}
\label{subsec:aor-methods-refpresb-cocomplete}

\Cref{thm:aor:refpresb-cocomplete} is the directed-colimit method used
for refinement chains.  A refinement chain is a directed diagram of
PreSB carriers whose arrows are annotated by the A1--A6 refinement
kinds.  The body argument separates two pieces of data: record and
scope components are collected by directed union, while carrier-level
data are read as compatible threads under the reverse projection maps.

The method is representative-object cocompleteness.  Rather than treating
the colimit as an unstructured quotient of all stages at once, the
construction chooses representative stage objects, carries the cocone data
through directed union, and tracks compatible carrier threads.  This is
enough for the paper's audit discipline: the body identifies which data must
be carried through a chain and how the representative witness sits in the
directed system.

The asymptotic-status arguments in \Cref{sec:aor-analysis} use this
method by asking what residual obligations remain at the limit of a
chain.  Finite cascade discharge is a statement about bounded cascade
depth at each stage; refinement-stable discharge also requires that
limit-stage residuals be paid by one of the four asymptotic statuses.
The representative-object construction is the bookkeeping device that
lets those two requirements be stated in the same scoped category.

\subsection{\texorpdfstring{Cost-table reproducibility and the
\texttt{definition\_metadata} discipline}{Cost-table reproducibility
and definition metadata discipline}}
\label{subsec:aor-methods-cost-table-reproducibility}

\Cref{def:aor:cost-table} introduces the base cost table used by the
cascade section.  The table is framework data: it specifies how each
residual type and base AOR-8 status generates downstream obligations.
It is not a theorem asserting that this table is uniquely forced by the
rest of the framework.  A different audited cost table would define a
different, though related, cascade discipline.

The reproducibility discipline is therefore explicit.  The base table
has \(12\times 11\) entries and belongs to the base AOR-8 status layer.
The representative non-trivial base rows include the
transport/\texttt{blocked}, source/\texttt{budgeted}, and
scale/\texttt{hidden} cases, together with structural defaults for the
remaining base rows.  The asymptotic limit rows, including the
contractive and asymptotic-budgeted limit cases, live in a separate
asymptotic table and are not counted as part of the 132-entry base
catalogue.

The body therefore treats the cost table as fixed framework data
rather than as a canonical or unique object.  Appendix E reserves
\Cref{tab:aor-cost-table-representative} for the representative
cost-table subset.  Downstream users of the AOR framework should cite
the cost table as part of the chosen audit discipline and should record any
altered cost rows explicitly.
 \section{Artifact source map}\label{app:artifact-map}

This appendix maps the AOR statements of record to their rendered body
locations and to any figure, table, or appendix support row.  It does
not duplicate the declaration data of \Cref{app:aor-formalization};
that material lives in the declaration maps of
\Cref{tab:paper4-lean-declaration-map-a,tab:paper4-lean-declaration-map-b}.

\begin{table}[!ht]
  \centering
  \scriptsize
  \setlength{\tabcolsep}{2.5pt}
  \renewcommand{\arraystretch}{1.08}
  \begin{tabularx}{\linewidth}{>{\raggedright\arraybackslash}p{0.045\linewidth} >{\raggedright\arraybackslash}p{0.28\linewidth} >{\raggedright\arraybackslash}p{0.20\linewidth} >{\raggedright\arraybackslash}p{0.25\linewidth} >{\raggedright\arraybackslash}X}
    \toprule
    \# & Paper label & Body section & Figure / table support & Appendix support \\
    \midrule
    01 & \path{def:aor:presb-carrier} & \Cref{sec:aor-problem-setup} & -- & \Cref{tab:paper4-lean-declaration-map-a} \\
    02 & \path{def:aor:scope-record} & \Cref{sec:aor-problem-setup} & -- & \Cref{tab:paper4-lean-declaration-map-a} \\
    03 & \path{def:aor:morphism-presb} & \Cref{sec:aor-problem-setup} & -- & \Cref{tab:paper4-lean-declaration-map-a} \\
    04 & \path{def:aor:aor-subcategory} & \Cref{sec:aor-problem-setup} & -- & \Cref{tab:paper4-lean-declaration-map-a} \\
    05 & \path{thm:aor:aor-full-subcat} & \Cref{sec:aor-problem-setup} & -- & \Cref{tab:paper4-lean-declaration-map-a} \\
    06 & \path{def:aor:cl-sb-operator} & \Cref{sec:aor-methods} & -- & \Cref{tab:paper4-lean-declaration-map-a} \\
    07 & \path{thm:aor:cl-sb-construction} & \Cref{sec:aor-methods} & -- & \Cref{tab:paper4-lean-declaration-map-a} \\
    08 & \path{thm:aor:cl-sb-universal-property} & \Cref{sec:aor-methods} & -- & \Cref{tab:paper4-lean-declaration-map-a} \\
    09 & \path{thm:aor:cl-sb-adjunction} & \Cref{sec:aor-methods} & -- & \Cref{tab:paper4-lean-declaration-map-a} \\
    10 & \path{def:aor:xi-sb} & \Cref{sec:aor-primary-results-defect} & -- & \Cref{tab:paper4-lean-declaration-map-a} \\
    11 & \path{def:aor:sigma-obs} & \Cref{sec:aor-primary-results-defect} & -- & \Cref{tab:paper4-lean-declaration-map-a} \\
    12 & \path{def:aor:sigma-stage} & \Cref{sec:aor-primary-results-defect} & -- & \Cref{tab:paper4-lean-declaration-map-a} \\
    13 & \path{def:aor:sigma-constraint} & \Cref{sec:aor-primary-results-defect} & -- & \Cref{tab:paper4-lean-declaration-map-a} \\
    14 & \path{def:aor:sigma-rewrite} & \Cref{sec:aor-primary-results-defect} & -- & \Cref{tab:paper4-lean-declaration-map-a} \\
    15 & \path{def:aor:sigma-holonomy} & \Cref{sec:aor-primary-results-defect} & -- & \Cref{tab:paper4-lean-declaration-map-a} \\
    16 & \path{def:aor:sigma-residual} & \Cref{sec:aor-primary-results-defect} & -- & \Cref{tab:paper4-lean-declaration-map-a} \\
    17 & \path{def:aor:sigma-nonclaim} & \Cref{sec:aor-primary-results-defect} & -- & \Cref{tab:paper4-lean-declaration-map-a} \\
    18 & \path{def:aor:sigma-meta} & \Cref{sec:aor-primary-results-defect} & -- & \Cref{tab:paper4-lean-declaration-map-a} \\
    19 & \path{thm:aor:eight-stratum-decomposition} & \Cref{sec:aor-primary-results-defect} & -- & \Cref{tab:paper4-lean-declaration-map-a} \\
    20 & \path{def:aor:residual-types-twelve} & \Cref{sec:aor-primary-results-atlas} & -- & \Cref{tab:paper4-lean-declaration-map-a} \\
    21 & \path{def:aor:residual-partial-order} & \Cref{sec:aor-primary-results-atlas} & -- & \Cref{tab:paper4-lean-declaration-map-a} \\
    22 & \path{def:aor:residual-atom} & \Cref{sec:aor-primary-results-atlas} & \Cref{tab:aor-forced-secondary-types} & \Cref{tab:paper4-lean-declaration-map-a} \\
    23 & \path{thm:aor:atom-decomposition} & \Cref{sec:aor-primary-results-atlas} & \Cref{tab:aor-forced-secondary-types} & \Cref{tab:paper4-lean-declaration-map-a} \\
    24 & \path{thm:aor:residual-product} & \Cref{sec:aor-primary-results-atlas} & \Cref{tab:aor-forced-secondary-types} & \Cref{tab:paper4-lean-declaration-map-a} \\
    \bottomrule
  \end{tabularx}
  \caption{AOR artifact source map, rows 1--24.}
  \label{tab:aor-artifact-map-a}
\end{table}

\begin{table}[!ht]
  \centering
  \scriptsize
  \setlength{\tabcolsep}{2.5pt}
  \renewcommand{\arraystretch}{1.08}
  \begin{tabularx}{\linewidth}{>{\raggedright\arraybackslash}p{0.045\linewidth} >{\raggedright\arraybackslash}p{0.28\linewidth} >{\raggedright\arraybackslash}p{0.20\linewidth} >{\raggedright\arraybackslash}p{0.25\linewidth} >{\raggedright\arraybackslash}X}
    \toprule
    \# & Paper label & Body section & Figure / table support & Appendix support \\
    \midrule
    25 & \path{thm:aor:forced-residual} & \Cref{sec:aor-primary-results-atlas} & \Cref{tab:aor-forced-secondary-types} & \Cref{tab:paper4-lean-declaration-map-a} \\
    26 & \path{def:aor:aor-8-statuses} & \Cref{sec:aor-primary-results-statuses} & -- & \Cref{tab:paper4-lean-declaration-map-a} \\
    27 & \path{def:aor:discharge-record} & \Cref{sec:aor-primary-results-statuses} & -- & \Cref{tab:paper4-lean-declaration-map-a} \\
    28 & \path{def:aor:discharge-witness} & \Cref{sec:aor-primary-results-statuses} & -- & \Cref{tab:paper4-lean-declaration-map-b} \\
    29 & \path{def:aor:cost-record} & \Cref{sec:aor-secondary-results} & \Cref{tab:aor-cost-table-representative} & \Cref{tab:paper4-lean-declaration-map-b} \\
    30 & \path{def:aor:down-operator} & \Cref{sec:aor-secondary-results} & \Cref{tab:aor-cost-table-representative} & \Cref{tab:paper4-lean-declaration-map-b} \\
    31 & \path{def:aor:cost-table} & \Cref{sec:aor-secondary-results} & \Cref{tab:aor-cost-table-representative} & \Cref{tab:paper4-lean-declaration-map-b} \\
    32 & \path{def:aor:cascade-graph} & \Cref{sec:aor-secondary-results} & -- & \Cref{tab:paper4-lean-declaration-map-b} \\
    33 & \path{thm:aor:cascade-bipartite} & \Cref{sec:aor-secondary-results} & -- & \Cref{tab:paper4-lean-declaration-map-b} \\
    34 & \path{thm:aor:cascade-termination} & \Cref{sec:aor-secondary-results} & -- & \Cref{tab:paper4-lean-declaration-map-b} \\
    35 & \path{thm:aor:cascade-confluence} & \Cref{sec:aor-secondary-results} & -- & \Cref{tab:paper4-lean-declaration-map-b,subsec:paper4-partial-disclosures} \\
    36 & \path{thm:aor:cross-cascade-composition} & \Cref{sec:aor-secondary-results} & -- & \Cref{tab:paper4-lean-declaration-map-b} \\
    37 & \path{def:aor:cascade-stable-aor} & \Cref{sec:aor-secondary-results} & \Cref{tab:aor-cost-table-representative} & \Cref{tab:paper4-lean-declaration-map-b} \\
    38 & \path{thm:aor:cascade-stable-characterisation} & \Cref{sec:aor-secondary-results} & \Cref{tab:aor-cost-table-representative} & \Cref{tab:paper4-lean-declaration-map-b} \\
    39 & \path{def:aor:refinement-morphism} & \Cref{sec:aor-analysis} & -- & \Cref{tab:paper4-lean-declaration-map-b} \\
    40 & \path{def:aor:refinement-kinds-six} & \Cref{sec:aor-analysis} & -- & \Cref{tab:paper4-lean-declaration-map-b} \\
    41 & \path{def:aor:refpresb-category} & \Cref{sec:aor-analysis} & -- & \Cref{tab:paper4-lean-declaration-map-b} \\
    42 & \path{thm:aor:refpresb-cocomplete} & \Cref{sec:aor-analysis} & -- & \Cref{tab:paper4-lean-declaration-map-b} \\
    43 & \path{def:aor:asymptotic-statuses} & \Cref{sec:aor-analysis} & -- & \Cref{tab:paper4-lean-declaration-map-b} \\
    44 & \path{thm:aor:asymptotic-soundness} & \Cref{sec:aor-analysis} & -- & \Cref{tab:paper4-lean-declaration-map-b} \\
    45 & \path{def:aor:refinement-stable-aor} & \Cref{sec:aor-analysis} & -- & \Cref{tab:paper4-lean-declaration-map-b} \\
    46 & \path{thm:aor:refinement-stable-characterisation} & \Cref{sec:aor-analysis} & -- & \Cref{tab:paper4-lean-declaration-map-b} \\
    47 & \path{thm:aor:hierarchy} & \Cref{sec:aor-analysis} & \Cref{tab:aor-counterexample-carriers} & \Cref{tab:paper4-lean-declaration-map-b,subsec:paper4-partial-disclosures} \\
    48 & \path{nonclaim:aor:no-strong-realism} & \Cref{sec:aor-discussion} & -- & -- \\
    \bottomrule
  \end{tabularx}
  \caption{AOR artifact source map, rows 25--48.}
  \label{tab:aor-artifact-map-b}
\end{table}

Rows 1--47 are the AOR statements of record and are accompanied by the
declaration maps in \Cref{app:aor-formalization}.  Row 48 is the
post-restructure nonclaim boundary carried in
\Cref{sec:aor-discussion}; it has no corresponding support row.
 \section{Lean Formalization Disclosure}
\label{app:aor-formalization}

\subsection{Trust base and methodology}
\label{subsec:paper4-trust-methodology}

The AOR Lean artifacts cited in this appendix were checked with Lean
4.28.0 \citep{deMouraUllrich2021Lean4}, whose dependent type theory
descends from Martin-L{\"o}f's intuitionistic type theory
\citep{MartinLof1984IntuitionisticTypeTheory} and the Calculus of
Constructions of \citet{CoquandHuet1988CalculusConstructions}.  The
project-wide trust base allowlisted in
\path{lean/manifests/trust_base.txt} comprises five core kernel
axioms: \texttt{propext}, \texttt{Classical.choice}, \texttt{Quot.sound},
\texttt{choice}, and \texttt{funext}.  The AOR-axis theorem closures
recorded by the manifest audit fit inside the smaller subset
$\{\texttt{propext},\ \texttt{Classical.choice},\ \texttt{Quot.sound}\}$.
The needles project remains mathlib-free \citep{MathlibCommunity2020Lean}
on this axis: categorical records, fixed-point packages, residual
types, cost rows, and refinement objects are represented locally by
Lean structures, inductive types, and \texttt{Prop}-valued predicates,
in the style familiar from other higher-order logic and dependent-type
proof assistants \citep{NipkowPaulsonWenzel2002IsabelleHOL}.
Project-local \texttt{sorry}, \texttt{axiom}, \texttt{opaque}, and
\texttt{constant} declarations are forbidden in the needles Lean
corpus; the manifest checker audits theorem closures against the
trust base.

\subsection{Mechanization disclosure tiers}
\label{subsec:paper4-mechanization-disclosure-tiers}

Body footnotes use a three-tier disclosure discipline keyed to the
alignment column of the per-label coverage table.

\begin{table}[H]
  \centering
  \small
  \begin{tabularx}{\linewidth}{>{\raggedright\arraybackslash}p{0.25\linewidth} >{\raggedright\arraybackslash}p{0.11\linewidth} >{\raggedright\arraybackslash}X}
    \toprule
    Alignment & Tier label & Body footnote wording template \\
    \midrule
    \texttt{faithful} & T1 & Verified in Lean as \(\langle\)declaration\(\rangle\); see \Cref{subsec:paper4-coverage-table}. \\
    \texttt{narrowed\_surrogate} & T2 & Mechanized in Lean as \(\langle\)declaration\(\rangle\); the Lean theorem proves \(\langle\)sub-claim\(\rangle\) under \(\langle\)narrowing\(\rangle\).  The full statement is proved here.  See \Cref{subsec:paper4-coverage-table}. \\
    \texttt{partial} with documented narrowing & T2-partial & Mechanized in Lean as \(\langle\)declaration\(\rangle\) as a partial mechanization: \(\langle\)recorded fragment\(\rangle\) is recorded; \(\langle\)deferred content\(\rangle\) is proved here.  See \Cref{subsec:paper4-coverage-table}. \\
    \bottomrule
  \end{tabularx}
  \caption{AOR body footnote disclosure tiers.}
  \label{tab:paper4-mechanization-disclosure-tiers}
\end{table}

This tier discipline is the canonical wording rule for body footnotes:
a reader can cross-check each theorem footnote against the alignment
column of the coverage table below.

\subsection{Standing narrowings}
\label{subsec:paper4-standing-narrowings}

Several narrowings apply uniformly across the AOR formalization.  First,
the \(\ClSB\) construction is presented in the paper as a
Knaster--Tarski monotone saturation on the complete lattice of
AOR-extending operations.  Lean records the complete-lattice side by a
\texttt{Prop}-valued placeholder
\(\texttt{completeLatticeOfOperations := True}\), while directly
exhibiting the operative closure package: idempotent closure, monotone
saturation, the universal lift, zero structural defect, and the
identity-field counit.

Second, the cost table is \texttt{definition\_metadata}.  The Lean
function \path{costTable} (in \path{SixBirdsNeedles.AOR.Cascades})
gives explicit non-trivial base rows for the pairs
\((\Dtransport,\blocked)\), \((\Dsource,\budgeted)\), and
\((\Dscale,\hidden)\).  The asymptotic limit rows live in
\path{asymptoticCostTable} (in \path{SixBirdsNeedles.AOR.Refinement})
and cover the \((\Dlimit,\contractive)\) and
\((\Dlimit,\asymptoticBudgeted)\) cases.  Remaining base rows use
structural defaults; \Cref{app:extended-material} reserves an appendix
table for a representative subset of the 132-entry catalogue.

Third, \texttt{closesLimitResidual} is no longer a uniform
\texttt{True} predicate.  The four asymptotic statuses
\(\eventualZero\), \(\contractive\), \(\summable\), and
\(\asymptoticBudgeted\) carry distinct content-bearing predicates in
\path{SixBirdsNeedles.AOR.Refinement}.

Fourth, \path{refpresbCocomplete} is a representative-object
narrowing.  Lean exhibits a colimit witness using the chain's base
object, with \texttt{True}-placeholder fields for the cocone,
directed-union, compatible-thread, functoriality, and
PreSB-typing-preservation clauses.  The existence of the witness is
captured; the full directed-colimit construction remains paper-side.

Fifth, the carrier-attached discharge assignment is represented by
the \path{AORCarrierWithDischarge} wrapper.  Cascade-stable AOR and
refinement-stable AOR are predicates over an AOR carrier paired with
a fixed residual-discharge assignment rather than over a free carrier
alone, ruling out the post-hoc choice of \texttt{zero} status for
every atom.  Any AOR-instance application, including the
NS-as-AOR-instance presentation in
\citet{Tsiokos2026NavierStokes}, inherits the wrapper.

\subsection{Per-label coverage table}
\label{subsec:paper4-coverage-table}

Tables~\ref{tab:paper4-coverage-a}--\ref{tab:paper4-coverage-c} cover all 47 AOR \texttt{mechanize\_now} rows.  The two partial rows (\path{thm:aor:cascade-confluence} and \path{thm:aor:hierarchy}) include their narrowing notes in the wording column.

\begin{landscape}

\begin{table}[!ht]
  \centering
  \scriptsize
  \setlength{\tabcolsep}{2pt}
  \renewcommand{\arraystretch}{1.0}
  \begin{tabularx}{\linewidth}{>{\raggedright\arraybackslash}p{0.04\linewidth} >{\raggedright\arraybackslash}p{0.22\linewidth} >{\raggedright\arraybackslash}p{0.11\linewidth} >{\raggedright\arraybackslash}p{0.11\linewidth} >{\raggedright\arraybackslash}p{0.20\linewidth} >{\raggedright\arraybackslash}X}
    \toprule
    \# & Paper label & Kind & Coverage & Alignment & Wording \\
    \midrule
    01 & \path{def:aor:presb-carrier} & \texttt{definition} & \texttt{definition} & \texttt{faithful} & Lean realizes the paper object as a Lean definition, structure, or inductive type. \\
    02 & \path{def:aor:scope-record} & \texttt{definition} & \texttt{definition} & \texttt{faithful} & Lean realizes the paper object as a Lean definition, structure, or inductive type. \\
    03 & \path{def:aor:morphism-presb} & \texttt{definition} & \texttt{definition} & \texttt{faithful} & Lean realizes the paper object as a Lean definition, structure, or inductive type. \\
    04 & \path{def:aor:aor-subcategory} & \texttt{definition} & \texttt{definition} & \texttt{faithful} & Lean realizes the paper object as a Lean definition, structure, or inductive type. \\
    05 & \path{thm:aor:aor-full-subcat} & \texttt{theorem} & \texttt{theorem} & \texttt{faithful} & Lean proves the stated structural claim. \\
    06 & \path{def:aor:cl-sb-operator} & \texttt{definition} & \texttt{definition} & \texttt{faithful} & Lean realizes the paper object as a Lean definition, structure, or inductive type. \\
    07 & \path{thm:aor:cl-sb-construction} & \texttt{theorem} & \texttt{theorem} & \texttt{narrowed\_surrogate} & Lean records a narrowed surrogate or audits the dependency package; the body carries the full mathematical proof. \\
    08 & \path{thm:aor:cl-sb-universal-property} & \texttt{theorem} & \texttt{theorem} & \texttt{narrowed\_surrogate} & Lean records a narrowed surrogate or audits the dependency package; the body carries the full mathematical proof. \\
    09 & \path{thm:aor:cl-sb-adjunction} & \texttt{theorem} & \texttt{theorem} & \texttt{narrowed\_surrogate} & Lean records a narrowed surrogate or audits the dependency package; the body carries the full mathematical proof. \\
    10 & \path{def:aor:xi-sb} & \texttt{definition} & \texttt{definition} & \texttt{faithful} & Lean realizes the paper object as a Lean definition, structure, or inductive type. \\
    11 & \path{def:aor:sigma-obs} & \texttt{definition} & \texttt{definition} & \texttt{faithful} & Lean realizes the paper object as a Lean definition, structure, or inductive type. \\
    12 & \path{def:aor:sigma-stage} & \texttt{definition} & \texttt{definition} & \texttt{faithful} & Lean realizes the paper object as a Lean definition, structure, or inductive type. \\
    13 & \path{def:aor:sigma-constraint} & \texttt{definition} & \texttt{definition} & \texttt{faithful} & Lean realizes the paper object as a Lean definition, structure, or inductive type. \\
    14 & \path{def:aor:sigma-rewrite} & \texttt{definition} & \texttt{definition} & \texttt{faithful} & Lean realizes the paper object as a Lean definition, structure, or inductive type. \\
    15 & \path{def:aor:sigma-holonomy} & \texttt{definition} & \texttt{definition} & \texttt{faithful} & Lean realizes the paper object as a Lean definition, structure, or inductive type. \\
    16 & \path{def:aor:sigma-residual} & \texttt{definition} & \texttt{definition} & \texttt{faithful} & Lean realizes the paper object as a Lean definition, structure, or inductive type. \\
    17 & \path{def:aor:sigma-nonclaim} & \texttt{definition} & \texttt{definition} & \texttt{faithful} & Lean realizes the paper object as a Lean definition, structure, or inductive type. \\
    18 & \path{def:aor:sigma-meta} & \texttt{definition} & \texttt{definition} & \texttt{faithful} & Lean realizes the paper object as a Lean definition, structure, or inductive type. \\
    \bottomrule
  \end{tabularx}
  \caption{AOR per-label coverage, rows 1--18.}
  \label{tab:paper4-coverage-a}
\end{table}

\begin{table}[!ht]
  \centering
  \scriptsize
  \setlength{\tabcolsep}{2pt}
  \renewcommand{\arraystretch}{1.0}
  \begin{tabularx}{\linewidth}{>{\raggedright\arraybackslash}p{0.04\linewidth} >{\raggedright\arraybackslash}p{0.22\linewidth} >{\raggedright\arraybackslash}p{0.11\linewidth} >{\raggedright\arraybackslash}p{0.11\linewidth} >{\raggedright\arraybackslash}p{0.20\linewidth} >{\raggedright\arraybackslash}X}
    \toprule
    \# & Paper label & Kind & Coverage & Alignment & Wording \\
    \midrule
    19 & \path{thm:aor:eight-stratum-decomposition} & \texttt{theorem} & \texttt{theorem} & \texttt{faithful} & Lean proves the stated structural claim. \\
    20 & \path{def:aor:residual-types-twelve} & \texttt{definition} & \texttt{definition} & \texttt{faithful} & Lean realizes the paper object as a Lean definition, structure, or inductive type. \\
    21 & \path{def:aor:residual-partial-order} & \texttt{definition} & \texttt{definition} & \texttt{faithful} & Lean realizes the paper object as a Lean definition, structure, or inductive type. \\
    22 & \path{def:aor:residual-atom} & \texttt{definition} & \texttt{definition} & \texttt{faithful} & Lean realizes the paper object as a Lean definition, structure, or inductive type. \\
    23 & \path{thm:aor:atom-decomposition} & \texttt{theorem} & \texttt{theorem} & \texttt{faithful} & Lean proves the stated structural claim. \\
    24 & \path{thm:aor:residual-product} & \texttt{theorem} & \texttt{theorem} & \texttt{faithful} & Lean proves the stated structural claim. \\
    25 & \path{thm:aor:forced-residual} & \texttt{theorem} & \texttt{theorem} & \texttt{faithful} & Lean proves the stated structural claim. \\
    26 & \path{def:aor:aor-8-statuses} & \texttt{definition} & \texttt{definition} & \texttt{faithful} & Lean realizes the paper object as a Lean definition, structure, or inductive type. \\
    27 & \path{def:aor:discharge-record} & \texttt{definition} & \texttt{definition} & \texttt{faithful} & Lean realizes the paper object as a Lean definition, structure, or inductive type. \\
    28 & \path{def:aor:discharge-witness} & \texttt{definition} & \texttt{definition} & \texttt{faithful} & Lean realizes the paper object as a Lean definition, structure, or inductive type. \\
    29 & \path{def:aor:cost-record} & \texttt{definition} & \texttt{definition} & \texttt{faithful} & Lean realizes the paper object as a Lean definition, structure, or inductive type. \\
    30 & \path{def:aor:down-operator} & \texttt{definition} & \texttt{definition} & \texttt{faithful} & Lean realizes the paper object as a Lean definition, structure, or inductive type. \\
    31 & \path{def:aor:cost-table} & \texttt{definition} & \texttt{definition} & \texttt{narrowed\_surrogate} & Lean records representative cost-table rows as definition metadata; the remaining rows are structural defaults. \\
    32 & \path{def:aor:cascade-graph} & \texttt{definition} & \texttt{definition} & \texttt{faithful} & Lean realizes the paper object as a Lean definition, structure, or inductive type. \\
    33 & \path{thm:aor:cascade-bipartite} & \texttt{theorem} & \texttt{theorem} & \texttt{narrowed\_surrogate} & Lean records a narrowed surrogate or audits the dependency package; the body carries the full mathematical proof. \\
    34 & \path{thm:aor:cascade-termination} & \texttt{theorem} & \texttt{theorem} & \texttt{narrowed\_surrogate} & Lean records a narrowed surrogate or audits the dependency package; the body carries the full mathematical proof. \\
    35 & \path{thm:aor:cascade-confluence} & \texttt{theorem} & \texttt{partial} & \texttt{narrowed\_surrogate} & Lean records the local confluence and well-founded discharge relation fragment; the body supplies the Newman implication and unique-normal-form conclusion as a classical input. \\
    36 & \path{thm:aor:cross-cascade-composition} & \texttt{theorem} & \texttt{theorem} & \texttt{faithful} & Lean proves the stated structural claim. \\
    \bottomrule
  \end{tabularx}
  \caption{AOR per-label coverage, rows 19--36.}
  \label{tab:paper4-coverage-b}
\end{table}

\begin{table}[!ht]
  \centering
  \scriptsize
  \setlength{\tabcolsep}{2pt}
  \renewcommand{\arraystretch}{1.0}
  \begin{tabularx}{\linewidth}{>{\raggedright\arraybackslash}p{0.04\linewidth} >{\raggedright\arraybackslash}p{0.22\linewidth} >{\raggedright\arraybackslash}p{0.11\linewidth} >{\raggedright\arraybackslash}p{0.11\linewidth} >{\raggedright\arraybackslash}p{0.20\linewidth} >{\raggedright\arraybackslash}X}
    \toprule
    \# & Paper label & Kind & Coverage & Alignment & Wording \\
    \midrule
    37 & \path{def:aor:cascade-stable-aor} & \texttt{definition} & \texttt{definition} & \texttt{faithful} & Lean realizes the paper object as a Lean definition, structure, or inductive type. \\
    38 & \path{thm:aor:cascade-stable-characterisation} & \texttt{theorem} & \texttt{theorem} & \texttt{narrowed\_surrogate} & Lean records a narrowed surrogate or audits the dependency package; the body carries the full mathematical proof. \\
    39 & \path{def:aor:refinement-morphism} & \texttt{definition} & \texttt{definition} & \texttt{faithful} & Lean realizes the paper object as a Lean definition, structure, or inductive type. \\
    40 & \path{def:aor:refinement-kinds-six} & \texttt{definition} & \texttt{definition} & \texttt{faithful} & Lean realizes the paper object as a Lean definition, structure, or inductive type. \\
    41 & \path{def:aor:refpresb-category} & \texttt{definition} & \texttt{definition} & \texttt{faithful} & Lean realizes the paper object as a Lean definition, structure, or inductive type. \\
    42 & \path{thm:aor:refpresb-cocomplete} & \texttt{theorem} & \texttt{theorem} & \texttt{narrowed\_surrogate} & Lean records a narrowed surrogate or audits the dependency package; the body carries the full mathematical proof. \\
    43 & \path{def:aor:asymptotic-statuses} & \texttt{definition} & \texttt{definition} & \texttt{faithful} & Lean realizes the paper object as a Lean definition, structure, or inductive type. \\
    44 & \path{thm:aor:asymptotic-soundness} & \texttt{theorem} & \texttt{theorem} & \texttt{narrowed\_surrogate} & Lean records a narrowed surrogate or audits the dependency package; the body carries the full mathematical proof. \\
    45 & \path{def:aor:refinement-stable-aor} & \texttt{definition} & \texttt{definition} & \texttt{faithful} & Lean realizes the paper object as a Lean definition, structure, or inductive type. \\
    46 & \path{thm:aor:refinement-stable-characterisation} & \texttt{theorem} & \texttt{theorem} & \texttt{narrowed\_surrogate} & Lean records a narrowed surrogate or audits the dependency package; the body carries the full mathematical proof. \\
    47 & \path{thm:aor:hierarchy} & \texttt{theorem} & \texttt{partial} & \texttt{narrowed\_surrogate} & Lean records the two witness carriers and their typed obstruction markers; full predicate-level non-membership verification remains framework future work. \\
    \bottomrule
  \end{tabularx}
  \caption{AOR per-label coverage, rows 37--47.}
  \label{tab:paper4-coverage-c}
\end{table}

\end{landscape}

\subsection{Two partial demotions: explicit narrowing disclosures}
\label{subsec:paper4-partial-disclosures}

Two partial demotions remain on the AOR axis; the third, \texttt{thm:aor:ns-refstableaor}, moved with the NS-as-AOR-instance content to \citet{Tsiokos2026NavierStokes} and is disclosed in that paper's formalization appendix.

\subsubsection{\texorpdfstring{\texttt{thm:aor:cascade-confluence}}{thm:aor:cascade-confluence}}

Lean mechanizes local confluence and the well-founded discharge relation.
The body supplies the Newman global-confluence implication and the
resulting unique-normal-form conclusion as a classical input.  The R1
repair removed the permissive \path{DischargeReachable.merge}
constructor, so the remaining local diamond is not obtained by an
arbitrary same-primary merge; it is paired with the finite termination
measure recorded by \path{cascadeTermination}.

\subsubsection{\texorpdfstring{\texttt{thm:aor:hierarchy}}{thm:aor:hierarchy}}

The hierarchy row is a witness-carrier theorem rather than a full
predicate-level classification theorem.  Lean exhibits the named finite carriers
\path{counterexampleAOROnly} and \path{counterexampleCascadeOnly},
together with typed obstruction inductives
\path{LimitAsymptoticNonpayment} and
\path{BlockedCascadeNontermination}.  Full predicate-level non-membership
verification remains framework future work and would require substantially
more of the 132-entry cost table than the present
\texttt{definition\_metadata} narrowing covers.

\subsection{Lean declaration map}
\label{subsec:paper4-declaration-map}

Tables~\ref{tab:paper4-lean-declaration-map-a} and \ref{tab:paper4-lean-declaration-map-b} map every mechanized AOR paper label to its Lean declaration path (47 rows total after the 2026-05-24 NS-AOR bridge restructure).

\begin{table}[H]
  \centering
  \scriptsize
  \setlength{\tabcolsep}{3pt}
  \renewcommand{\arraystretch}{1.0}
  \begin{tabularx}{\linewidth}{>{\raggedright\arraybackslash}p{0.06\linewidth} >{\raggedright\arraybackslash}p{0.35\linewidth} >{\raggedright\arraybackslash}X}
    \toprule
    \# & Paper label & Lean declaration path \\
    \midrule
    01 & \path{def:aor:presb-carrier} & \path{SixBirdsNeedles.AOR.PreSB.preSBCarrier} \\
    02 & \path{def:aor:scope-record} & \path{SixBirdsNeedles.AOR.PreSB.scopeRecord} \\
    03 & \path{def:aor:morphism-presb} & \path{SixBirdsNeedles.AOR.PreSB.morphismPreSB} \\
    04 & \path{def:aor:aor-subcategory} & \path{SixBirdsNeedles.AOR.PreSB.aorSubcategory} \\
    05 & \path{thm:aor:aor-full-subcat} & \path{SixBirdsNeedles.AOR.PreSB.aorFullSubcat} \\
    06 & \path{def:aor:cl-sb-operator} & \path{SixBirdsNeedles.AOR.ClSB.clSBOperator} \\
    07 & \path{thm:aor:cl-sb-construction} & \path{SixBirdsNeedles.AOR.ClSB.clSBConstruction} \\
    08 & \path{thm:aor:cl-sb-universal-property} & \path{SixBirdsNeedles.AOR.ClSB.clSBUniversalProperty} \\
    09 & \path{thm:aor:cl-sb-adjunction} & \path{SixBirdsNeedles.AOR.ClSB.clSBAdjunction} \\
    10 & \path{def:aor:xi-sb} & \path{SixBirdsNeedles.AOR.DefectStrata.xiSB} \\
    11 & \path{def:aor:sigma-obs} & \path{SixBirdsNeedles.AOR.DefectStrata.sigmaObs} \\
    12 & \path{def:aor:sigma-stage} & \path{SixBirdsNeedles.AOR.DefectStrata.sigmaStage} \\
    13 & \path{def:aor:sigma-constraint} & \path{SixBirdsNeedles.AOR.DefectStrata.sigmaConstraint} \\
    14 & \path{def:aor:sigma-rewrite} & \path{SixBirdsNeedles.AOR.DefectStrata.sigmaRewrite} \\
    15 & \path{def:aor:sigma-holonomy} & \path{SixBirdsNeedles.AOR.DefectStrata.sigmaHolonomy} \\
    16 & \path{def:aor:sigma-residual} & \path{SixBirdsNeedles.AOR.DefectStrata.sigmaResidual} \\
    17 & \path{def:aor:sigma-nonclaim} & \path{SixBirdsNeedles.AOR.DefectStrata.sigmaNonclaim} \\
    18 & \path{def:aor:sigma-meta} & \path{SixBirdsNeedles.AOR.DefectStrata.sigmaMeta} \\
    19 & \path{thm:aor:eight-stratum-decomposition} & \path{SixBirdsNeedles.AOR.DefectStrata.eightStratumDecomposition} \\
    20 & \path{def:aor:residual-types-twelve} & \path{SixBirdsNeedles.AOR.ResidualAtlas.residualTypesTwelve} \\
    21 & \path{def:aor:residual-partial-order} & \path{SixBirdsNeedles.AOR.ResidualAtlas.residualPartialOrder} \\
    22 & \path{def:aor:residual-atom} & \path{SixBirdsNeedles.AOR.ResidualAtlas.residualAtom} \\
    23 & \path{thm:aor:atom-decomposition} & \path{SixBirdsNeedles.AOR.ResidualAtlas.atomDecomposition} \\
    24 & \path{thm:aor:residual-product} & \path{SixBirdsNeedles.AOR.ResidualAtlas.residualProduct} \\
    25 & \path{thm:aor:forced-residual} & \path{SixBirdsNeedles.AOR.ResidualAtlas.forcedResidual} \\
    26 & \path{def:aor:aor-8-statuses} & \path{SixBirdsNeedles.AOR.Cascades.AOR8Status} \\
    27 & \path{def:aor:discharge-record} & \path{SixBirdsNeedles.AOR.Cascades.DischargeWitness} \\
    \bottomrule
  \end{tabularx}
  \caption{AOR Lean declaration map, rows 1--27.}
  \label{tab:paper4-lean-declaration-map-a}
\end{table}

\begin{table}[H]
  \centering
  \scriptsize
  \setlength{\tabcolsep}{3pt}
  \renewcommand{\arraystretch}{1.0}
  \begin{tabularx}{\linewidth}{>{\raggedright\arraybackslash}p{0.06\linewidth} >{\raggedright\arraybackslash}p{0.35\linewidth} >{\raggedright\arraybackslash}X}
    \toprule
    \# & Paper label & Lean declaration path \\
    \midrule
    28 & \path{def:aor:discharge-witness} & \path{SixBirdsNeedles.AOR.Cascades.DischargeWitnessFamily} \\
    29 & \path{def:aor:cost-record} & \path{SixBirdsNeedles.AOR.Cascades.costRecord} \\
    30 & \path{def:aor:down-operator} & \path{SixBirdsNeedles.AOR.Cascades.downOperator} \\
    31 & \path{def:aor:cost-table} & \path{SixBirdsNeedles.AOR.Cascades.costTable} \\
    32 & \path{def:aor:cascade-graph} & \path{SixBirdsNeedles.AOR.Cascades.cascadeGraph} \\
    33 & \path{thm:aor:cascade-bipartite} & \path{SixBirdsNeedles.AOR.Cascades.cascadeBipartite} \\
    34 & \path{thm:aor:cascade-termination} & \path{SixBirdsNeedles.AOR.Cascades.cascadeTermination} \\
    35 & \path{thm:aor:cascade-confluence} & \path{SixBirdsNeedles.AOR.Cascades.cascadeConfluence} \\
    36 & \path{thm:aor:cross-cascade-composition} & \path{SixBirdsNeedles.AOR.Cascades.crossCascadeComposition} \\
    37 & \path{def:aor:cascade-stable-aor} & \path{SixBirdsNeedles.AOR.Cascades.cascadeStableAOR} \\
    38 & \path{thm:aor:cascade-stable-characterisation} & \path{SixBirdsNeedles.AOR.Cascades.cascadeStableCharacterisation} \\
    39 & \path{def:aor:refinement-morphism} & \path{SixBirdsNeedles.AOR.Refinement.refinementMorphism} \\
    40 & \path{def:aor:refinement-kinds-six} & \path{SixBirdsNeedles.AOR.Refinement.refinementKindsSix} \\
    41 & \path{def:aor:refpresb-category} & \path{SixBirdsNeedles.AOR.Refinement.refpresbCategory} \\
    42 & \path{thm:aor:refpresb-cocomplete} & \path{SixBirdsNeedles.AOR.Refinement.refpresbCocomplete} \\
    43 & \path{def:aor:asymptotic-statuses} & \path{SixBirdsNeedles.AOR.Refinement.asymptoticStatuses} \\
    44 & \path{thm:aor:asymptotic-soundness} & \path{SixBirdsNeedles.AOR.Refinement.asymptoticSoundness} \\
    45 & \path{def:aor:refinement-stable-aor} & \path{SixBirdsNeedles.AOR.Refinement.refinementStableAOR} \\
    46 & \path{thm:aor:refinement-stable-characterisation} & \path{SixBirdsNeedles.AOR.Refinement.refinementStableCharacterisation} \\
    47 & \path{thm:aor:hierarchy} & \path{SixBirdsNeedles.AOR.Refinement.hierarchy} \\
    \bottomrule
  \end{tabularx}
  \caption{AOR Lean declaration map, rows 28--47.}
  \label{tab:paper4-lean-declaration-map-b}
\end{table}

\subsection{Scope of this appendix}
\label{subsec:paper4-restructure-note}

This appendix covers the 47 mechanized AOR rows listed in
\path{lean/manifests/aor_manifest.toml}.  The NS-as-AOR-instance
content and its formalization disclosure, including the
\texttt{R4SourceAdmission} carrier-space field extension, are carried
by \citet{Tsiokos2026NavierStokes}.
 \section{Extended tables and nonclaims ledger}\label{app:extended-material}

This appendix collects framework data that are too detailed for the
body: the representative cost-table subset, the hierarchy
witness carriers, the forced-secondary table, and the complete
scope ledger.  The tables are audit data, not additional theorem
statements.

\subsection{Representative cost-table subset}
\label{subsec:aor-cost-table-representative}

\Cref{def:aor:cost-table} introduces the base \(132\)-entry
\(\mathrm{Cost}\) table as framework data.  The table below prints a
representative \(22\)-row base subset plus the two asymptotic rows used
for limit-stage refinement discharge.  The base rows and the asymptotic
rows are deliberately separated in the notes column: the latter belong
to the refinement-chain asymptotic table, not to the \(132\)-entry base
count.

\begin{table}[!ht]
  \centering
  \scriptsize
  \setlength{\tabcolsep}{3pt}
  \renewcommand{\arraystretch}{1.06}
  \begin{tabularx}{\linewidth}{>{\raggedright\arraybackslash}p{0.18\linewidth} >{\raggedright\arraybackslash}p{0.18\linewidth} >{\raggedright\arraybackslash}X >{\raggedright\arraybackslash}p{0.24\linewidth}}
    \toprule
    Primary residual type & Status & Downstream multiset & Reproducibility note \\
    \midrule
    \(\Dtransport\) & \texttt{blocked} & \(\Dtransport\) at \texttt{budgeted}; \(\Drole\) at \texttt{zero} & Explicit non-trivial base row. \\
    \(\Dsource\) & \texttt{budgeted} & \(\Dtransport\), \(\Dscale\), \(\Dtarget\), \(\Drole\) at \texttt{zero} & Explicit non-trivial base row. \\
    \(\Dscale\) & \texttt{hidden} & \(\Dscale\) at \texttt{blocked} & Explicit non-trivial base row. \\
    any base type & \texttt{zero} & \(\varnothing\) & Cost-free base row. \\
    any base type & \texttt{approved\_other} & \(\varnothing\) & Cost-free base row. \\
    \(\Ddescent\) & other base statuses & \(\Ddescent,\Dtarget\) at \texttt{zero} & Structural default. \\
    \(\Dtransport\) & other base statuses & \(\Dtransport\), \(\Dscale\), \(\Dlimit\), \(\Drole\), \(\Dtarget\), \(\Dpresentation\) at \texttt{zero} & Structural default. \\
    \(\Dinterface\) & other base statuses & \(\Dinterface\), \(\Dlocal\), \(\Dglobal\) at \texttt{zero} & Structural default. \\
    \(\Dlocal\) & other base statuses & \(\Dlocal\), \(\Dglobal\) at \texttt{zero} & Structural default. \\
    \(\Dglobal\) & other base statuses & \(\Dglobal\), \(\Dlocal\) at \texttt{zero} & Structural default. \\
    \(\Dpresentation\) & other base statuses & \(\Dpresentation\) at \texttt{zero} & Structural default. \\
    \(\Dsource\) & other base statuses & \(\Dsource\), \(\Dtransport\), \(\Dscale\), \(\Drole\), \(\Dtarget\) at \texttt{zero} & Structural default. \\
    \(\Dhorizon\) & other base statuses & \(\Dhorizon\) at \texttt{zero} & Structural default. \\
    \(\Dscale\) & other base statuses & \(\Dscale\), \(\Dlimit\), \(\Dhorizon\), \(\Dsource\), \(\Dtarget\), \(\Dpresentation\) at \texttt{zero} & Structural default. \\
    \(\Dlimit\) & other base statuses & \(\Dlimit\), \(\Dscale\), \(\Dhorizon\) at \texttt{zero} & Structural default. \\
    \(\Dsym\) & other base statuses & \(\Dsym\) at \texttt{zero} & Structural default. \\
    \(\Drole\) & other base statuses & \(\Drole\), \(\Dtarget\) at \texttt{zero} & Structural default. \\
    \(\Dtarget\) & other base statuses & \(\Dtarget\), \(\Drole\), \(\Dpresentation\) at \texttt{zero} & Structural default. \\
    any type except explicit cases & \texttt{budgeted} & type default at \texttt{zero} & Status-family representative. \\
    any type except explicit cases & \texttt{bridged}, \texttt{moduli}, \texttt{selected} & type default at \texttt{zero} & Status-family representative. \\
    any type except explicit cases & \texttt{probabilistic}, \texttt{out\_of\_scope}, \texttt{nonclaim} & type default at \texttt{zero} & Status-family representative. \\
    any type except explicit cases & \texttt{blocked} or \texttt{hidden} & type default at \texttt{zero} & Status-family representative. \\
    \(\Dlimit\) & \texttt{contractive} & \(\varnothing\) & Asymptotic companion row. \\
    \(\Dlimit\) & \texttt{asymptotic\_budgeted} & \(\Dlimit\) at \texttt{summable} & Asymptotic companion row. \\
    \bottomrule
  \end{tabularx}
  \caption{Representative cost-table subset: 22 grouped base rows plus two asymptotic companion rows.}
  \label{tab:aor-cost-table-representative}
\end{table}

\subsection{Witness carriers for the hierarchy}
\label{subsec:aor-counterexample-carriers}

\Cref{thm:aor:hierarchy} uses named finite carriers to mark the witness
boundaries in the hierarchy.  The table records the construction data:
carrier name, scope, residual type, status pattern, and typed obstruction
marker.  Full predicate-level verification for the hierarchy witnesses
remains framework future work.

\begin{landscape}

\begin{table}[!ht]
  \centering
  \small
  \setlength{\tabcolsep}{6pt}
  \renewcommand{\arraystretch}{1.12}
  \begin{tabularx}{\linewidth}{>{\raggedright\arraybackslash}p{0.24\linewidth} >{\raggedright\arraybackslash}p{0.10\linewidth} >{\raggedright\arraybackslash}p{0.08\linewidth} >{\raggedright\arraybackslash}X >{\raggedright\arraybackslash}p{0.24\linewidth}}
    \toprule
    Carrier name & Scope & Residual type & Assignment / status pattern & Obstruction marker \\
    \midrule
    \path{counterexampleAOROnly} & finite toy scope & \(\Dtransport\) & declared transport residual assigned \texttt{blocked}, producing a non-terminating obstruction pattern in the chosen cost-table model & \path{BlockedCascadeNontermination} \\
    \path{counterexampleCascadeOnly} & finite toy scope & \(\Dlimit\) & finite cascade closes, but the surviving limit residual is not paid by the asymptotic status required for refinement stability & \path{LimitAsymptoticNonpayment} \\
    \bottomrule
  \end{tabularx}
  \caption{Named witness carriers used by the hierarchy theorem.}
  \label{tab:aor-counterexample-carriers}
\end{table}

\end{landscape}

\subsection{Forced-secondary-type table}
\label{subsec:aor-forced-secondary-types}

The forced-secondary rule in \Cref{thm:aor:forced-residual} prevents a
discharge from closing only the headline residual while leaving its
typed dependencies open.  The local and global rows are grouped because
they form a reciprocal pair.

\begin{table}[!ht]
  \centering
  \small
  \setlength{\tabcolsep}{5pt}
  \renewcommand{\arraystretch}{1.1}
  \begin{tabularx}{\linewidth}{>{\raggedright\arraybackslash}p{0.28\linewidth} >{\raggedright\arraybackslash}X}
    \toprule
    Primary residual type & Forced secondary types \\
    \midrule
    \(\Dsource\) & \(\Dtransport,\Dscale,\Drole,\Dtarget\) \\
    \(\Dtransport\) & \(\Dscale,\Dlimit,\Drole,\Dtarget,\Dpresentation\) \\
    \(\Dscale\) & \(\Dlimit,\Dhorizon,\Dsource,\Dtarget,\Dpresentation\) \\
    \(\Dlimit\) & \(\Dscale,\Dhorizon\) \\
    \(\Dtarget\) & \(\Drole,\Dpresentation\) \\
    \(\Drole\) & \(\Dtarget\) \\
    \(\Ddescent\) & \(\Dtarget\) \\
    \(\Dinterface\) & \(\Dlocal,\Dglobal\) \\
    \(\Dpresentation\) & none \\
    \(\Dhorizon\) & none \\
    \(\Dsym\) & none \\
    \(\Dlocal,\Dglobal\) & reciprocal: \(\Dlocal\) forces \(\Dglobal\), and \(\Dglobal\) forces \(\Dlocal\) \\
    \bottomrule
  \end{tabularx}
  \caption{Forced secondary residual types used by \Cref{thm:aor:forced-residual}.}
  \label{tab:aor-forced-secondary-types}
\end{table}

\subsection{Complete AOR nonclaims ledger}
\label{subsec:aor-nonclaims-ledger}

\begin{description}[style=nextline]
\item[No strong-realism reading.]
Membership in \(\AOR_{\mathfrak S}\) is a record-level audit predicate
inside a declared scope, not a metaphysical claim about the carrier's
underlying domain.  This is the canonical nonclaim of
\Cref{nonclaim:aor:no-strong-realism}.

\item[No applications outside the audit-closed-carrier discipline.]
The framework stops at the AOR meta-theory.  Downstream applications
must supply their own scope record, residual atlas, cost table, and
discharge assignment; the current cited consumer is
\citet{Tsiokos2026NavierStokes}.

\item[No canonical-cost-table claim.]
The base cost table and the asymptotic table are part of the chosen
framework data.  They are not uniqueness theorems, and alternative
audited row choices would define related but different cascade systems.

\item[Hierarchy witness verification remains future work.]
The hierarchy theorem exhibits named carriers and typed obstruction
markers, but full predicate-level verification for the hierarchy witnesses
remains future work pending a fuller cost-table enumeration.

\item[Classical confluence input.]
The cascade-confluence theorem proves the local diamond and well-founded
termination, then applies Newman's lemma to obtain global confluence and
unique normal form.
\end{description}
 
\clearpage
\bibliographystyle{plainnat}
\bibliography{references}

\end{document}
